% main.tex -- OLS Prime paper, first draft (2026-11-19).
% Brief: .claude/briefs/iteration-5/paper-draft-1.md, under paper.md.
% Every number is a macro of numbers.tex, generated by make_numbers.py from the line its comment
% names. Wave 4's rates, H15's verdict and the case study's analyses wait for audit phase 2a and
% appear as \todo{after audit 2a}; their tables are laid out with empty cells.
% Typed by hand, as identifiers rather than measurements: dates, months and years of the periods,
% commit hashes, arXiv identifiers, hypothesis and section numbers, clause numbers.
%
% Build (MiKTeX, nothing installed):
%   cd python && uv run --no-sync python ../paper/make_numbers.py
%   cd ../paper && pdflatex -interaction=nonstopmode --disable-installer main.tex
%   bibtex main && pdflatex ... main.tex && pdflatex ... main.tex

\documentclass[11pt]{article}
\usepackage[T1]{fontenc}
\usepackage[utf8]{inputenc}
\usepackage{lmodern}
\usepackage{textcomp}
\usepackage[margin=1in]{geometry}
\usepackage{amsmath,amssymb}
\usepackage{booktabs}
\usepackage{array}
\usepackage{tabularx}
\usepackage{multirow}
\usepackage{xcolor}
\usepackage{enumitem}
\usepackage{caption}
\usepackage{url}
\usepackage[numbers,sort&compress]{natbib}
\usepackage{microtype}
\usepackage[hidelinks]{hyperref}
\usepackage{cleveref}

% numbers.tex -- GENERATED by paper/make_numbers.py; do not edit by hand.
% Every macro below is read from the line named in its comment.
% Wave 4's rates, H15's verdict and the case study's analyses are not here: they wait
% for audit phase 2a and appear in main.tex as \todo{after audit 2a}.

% docs/RESULTS.md:8, row (2026-10-28 | H1)
%   pattern: (\d+) honest models, (\d+) controls
%   matched: 11 honest models, 27 controls
\newcommand{\NumHoneHonest}{11}  % raw: 11
\newcommand{\NumHoneControls}{27}  % raw: 27

% docs/RESULTS.md:8, row (2026-10-28 | H1)
%   pattern: 1d: (\d+) cuts per pair
%   matched: 1d: 8 cuts per pair
\newcommand{\NumHonePairCuts}{8}  % raw: 8

% docs/RESULTS.md:8, row (2026-10-28 | H1)
%   pattern: multiplied by (\d+) in one run
%   matched: multiplied by 3 in one run
\newcommand{\NumHoneFactor}{3}  % raw: 3

% docs/RESULTS.md:12, row (2026-10-28 | H1 (second run; erratum to the first H1 row))
%   pattern: cheater is caught at (\d+) of (\d+) cuts
%   matched: cheater is caught at 8 of 8 cuts
\newcommand{\NumHoneCheatCut}{8}  % raw: 8
\newcommand{\NumHoneCheatCuts}{8}  % raw: 8

% docs/RESULTS.md:12, row (2026-10-28 | H1 (second run; erratum to the first H1 row))
%   pattern: smallest catch is ([\d.]+e-\d+) on BTCUSDT and ([\d.]+e-\d+) on ETHUSDT
%   matched: smallest catch is 7.1e-4 on BTCUSDT and 7.3e-4 on ETHUSDT
\newcommand{\NumHoneSmallBTC}{\ensuremath{7.1\times 10^{-4}}}  % raw: 7.1e-4
\newcommand{\NumHoneSmallETH}{\ensuremath{7.3\times 10^{-4}}}  % raw: 7.3e-4

% docs/RESULTS.md:13, row (2026-10-29 | H1 (panel))
%   pattern: arms on the (\d+)-member sub-panel
%   matched: arms on the 10-member sub-panel
\newcommand{\NumHonePanelMembers}{10}  % raw: 10

% docs/RESULTS.md:13, row (2026-10-29 | H1 (panel))
%   pattern: hourly bars: (\d+) cuts
%   matched: hourly bars: 8 cuts
\newcommand{\NumHonePanelCuts}{8}  % raw: 8

% docs/RESULTS.md:14, row (2026-10-29 | H1 (adversary, descriptive))
%   pattern: (\d{1,3}(?:,\d{3})+|\d+) random leaks \(seed (\d+)\)
%   matched: 10,000 random leaks (seed 20261029)
\newcommand{\NumAdvN}{10,000}  % raw: 10,000
\newcommand{\NumAdvSeed}{20261029}  % raw: 20261029

% docs/RESULTS.md:14, row (2026-10-29 | H1 (adversary, descriptive))
%   pattern: at H1's (\d+) cuts and a long schedule of (\d+) cuts
%   matched: at H1's 16 cuts and a long schedule of 128 cuts
\newcommand{\NumAdvCuts}{16}  % raw: 16
\newcommand{\NumAdvLongCuts}{128}  % raw: 128

% docs/RESULTS.md:14, row (2026-10-29 | H1 (adversary, descriptive))
%   pattern: Caught at one of H1's cuts at least: (\d{1,3}(?:,\d{3})+|\d+) of (\d{1,3}(?:,\d{3})+|\d+), ([-+<U+2212>]?(?:\d{1,3}(?:,\d{3})+|\d+)(?:\.\d+)?) %, 95 % Wilson \[([-+<U+2212>]?(?:\d{1,3}(?:,\d{3})+|\d+)(?:\.\d+)?) %, ([-+<U+2212>]?(?:\d{1,3}(?:,\d{3})+|\d+)(?:\.\d+)?) %\]
%   matched: Caught at one of H1's cuts at least: 9,795 of 10,000, 98.0 %, 95 % Wilson [97.7 %, 98.2 %]
\newcommand{\NumAdvCaught}{9,795}  % raw: 9,795
\newcommand{\NumAdvTotal}{10,000}  % raw: 10,000
\newcommand{\NumAdvRate}{98.0}  % raw: 98.0
\newcommand{\NumAdvRateLo}{97.7}  % raw: 97.7
\newcommand{\NumAdvRateHi}{98.2}  % raw: 98.2

% docs/RESULTS.md:14, row (2026-10-29 | H1 (adversary, descriptive))
%   pattern: Mean per-cut rate at H1's cuts: ([-+<U+2212>]?(?:\d{1,3}(?:,\d{3})+|\d+)(?:\.\d+)?) %, 95 % \[([-+<U+2212>]?(?:\d{1,3}(?:,\d{3})+|\d+)(?:\.\d+)?) %, ([-+<U+2212>]?(?:\d{1,3}(?:,\d{3})+|\d+)(?:\.\d+)?) %\]
%   matched: Mean per-cut rate at H1's cuts: 86.9 %, 95 % [86.4 %, 87.3 %]
\newcommand{\NumAdvPerCut}{86.9}  % raw: 86.9
\newcommand{\NumAdvPerCutLo}{86.4}  % raw: 86.4
\newcommand{\NumAdvPerCutHi}{87.3}  % raw: 87.3

% docs/RESULTS.md:14, row (2026-10-29 | H1 (adversary, descriptive))
%   pattern: Honest twins caught: (\d+) of (\d{1,3}(?:,\d{3})+|\d+)
%   matched: Honest twins caught: 0 of 10,000
\newcommand{\NumAdvTwinsCaught}{0}  % raw: 0
\newcommand{\NumAdvTwinsTotal}{10,000}  % raw: 10,000

% docs/RESULTS.md:14, row (2026-10-29 | H1 (adversary, descriptive))
%   pattern: (\d{1,3}(?:,\d{3})+|\d+) of them \(([-+<U+2212>]?(?:\d{1,3}(?:,\d{3})+|\d+)(?:\.\d+)?) %\) are constant at the compared rows and cannot
%   matched: 1,166 of them (11.7 %) are constant at the compared rows and cannot
\newcommand{\NumAdvTwinsConst}{1,166}  % raw: 1,166
\newcommand{\NumAdvTwinsConstPct}{11.7}  % raw: 11.7

% docs/RESULTS.md:14, row (2026-10-29 | H1 (adversary, descriptive))
%   pattern: Designed leaks: (\d{1,3}(?:,\d{3})+|\d+) of (\d{1,3}(?:,\d{3})+|\d+) caught
%   matched: Designed leaks: 1,247 of 1,247 caught
\newcommand{\NumAdvDesigned}{1,247}  % raw: 1,247
\newcommand{\NumAdvDesignedTotal}{1,247}  % raw: 1,247

% docs/RESULTS.md:14, row (2026-10-29 | H1 (adversary, descriptive))
%   pattern: eps (1e-10) ([-+<U+2212>]?(?:\d{1,3}(?:,\d{3})+|\d+)(?:\.\d+)?) %
%   matched: eps 1e-10 2.4 %
\newcommand{\NumAdvEpsMin}{\ensuremath{10^{-10}}}  % raw: 1e-10
\newcommand{\NumAdvEpsMinRate}{2.4}  % raw: 2.4

% docs/RESULTS.md:14, row (2026-10-29 | H1 (adversary, descriptive))
%   pattern: (\d+) leaks are caught at no cut of H1; (\d+) of them are constant
%   matched: 205 leaks are caught at no cut of H1; 163 of them are constant
\newcommand{\NumAdvMissed}{205}  % raw: 205
\newcommand{\NumAdvMissedConst}{163}  % raw: 163

% docs/RESULTS.md:22, row (2026-10-30 | H1 (blind spots, descriptive))
%   pattern: K = (\d+) twins per cut
%   matched: K = 8 twins per cut
\newcommand{\NumTwinsK}{8}  % raw: 8

% docs/RESULTS.md:22, row (2026-10-30 | H1 (blind spots, descriptive))
%   pattern: constant at the compared rows (\d+) \(([-+<U+2212>]?(?:\d{1,3}(?:,\d{3})+|\d+)(?:\.\d+)?) %, Wilson \[([-+<U+2212>]?(?:\d{1,3}(?:,\d{3})+|\d+)(?:\.\d+)?), ([-+<U+2212>]?(?:\d{1,3}(?:,\d{3})+|\d+)(?:\.\d+)?)\]\)
%   matched: constant at the compared rows 193 (94.1 %, Wilson [90.0, 96.6])
\newcommand{\NumTwinsConst}{193}  % raw: 193
\newcommand{\NumTwinsConstPct}{94.1}  % raw: 94.1
\newcommand{\NumTwinsConstLo}{90.0}  % raw: 90.0
\newcommand{\NumTwinsConstHi}{96.6}  % raw: 96.6

% docs/RESULTS.md:22, row (2026-10-30 | H1 (blind spots, descriptive))
%   pattern: unchanged under every twin (\d+) \(([-+<U+2212>]?(?:\d{1,3}(?:,\d{3})+|\d+)(?:\.\d+)?) %
%   matched: unchanged under every twin 5 (2.4 %
\newcommand{\NumTwinsUnch}{5}  % raw: 5
\newcommand{\NumTwinsUnchPct}{2.4}  % raw: 2.4

% docs/RESULTS.md:22, row (2026-10-30 | H1 (blind spots, descriptive))
%   pattern: true blind spots (\d+) \(([-+<U+2212>]?(?:\d{1,3}(?:,\d{3})+|\d+)(?:\.\d+)?) %
%   matched: true blind spots 7 (3.4 %
\newcommand{\NumTwinsBlind}{7}  % raw: 7
\newcommand{\NumTwinsBlindPct}{3.4}  % raw: 3.4

% docs/RESULTS.md:22, row (2026-10-30 | H1 (blind spots, descriptive))
%   pattern: Caught leaks moved by some twin: (\d{1,3}(?:,\d{3})+|\d+) of (\d{1,3}(?:,\d{3})+|\d+), ([-+<U+2212>]?(?:\d{1,3}(?:,\d{3})+|\d+)(?:\.\d+)?) %
%   matched: Caught leaks moved by some twin: 9,626 of 9,795, 98.3 %
\newcommand{\NumTwinsMoved}{9,626}  % raw: 9,626
\newcommand{\NumTwinsMovedOf}{9,795}  % raw: 9,795
\newcommand{\NumTwinsMovedPct}{98.3}  % raw: 98.3

% docs/RESULTS.md:32, row (2026-11-18 | H1 (correction of the rows of 2026-10-29 and 2026-10-30, descriptive))
%   pattern: (\d+) misses at H1's cuts: both schedules' flag (\d+) constant / (\d+) unchanged / (\d+) blind spots; H1's cuts only (\d+) / (\d+) / (\d+); nested, (\d+) between them
%   matched: 205 misses at H1's cuts: both schedules' flag 163 constant / 22 unchanged / 20 blind spots; H1's cuts only 193 / 5 / 7; nested, 30 between them
\newcommand{\NumCorrMisses}{205}  % raw: 205
\newcommand{\NumCorrBothConst}{163}  % raw: 163
\newcommand{\NumCorrBothUnch}{22}  % raw: 22
\newcommand{\NumCorrBothBlind}{20}  % raw: 20
\newcommand{\NumCorrHoneConst}{193}  % raw: 193
\newcommand{\NumCorrHoneUnch}{5}  % raw: 5
\newcommand{\NumCorrHoneBlind}{7}  % raw: 7
\newcommand{\NumCorrBetween}{30}  % raw: 30

% docs/RESULTS.md:32, row (2026-11-18 | H1 (correction of the rows of 2026-10-29 and 2026-10-30, descriptive))
%   pattern: pools all (\d+) cuts of its two schedules
%   matched: pools all 144 cuts of its two schedules
\newcommand{\NumCorrAllCuts}{144}  % raw: 144

% docs/RESULTS.md:25, row (2026-10-31 | H1 (encoder arms))
%   pattern: Brier skill score ([-+<U+2212>]?(?:\d{1,3}(?:,\d{3})+|\d+)(?:\.\d+)?) \[([-+<U+2212>]?(?:\d{1,3}(?:,\d{3})+|\d+)(?:\.\d+)?), ([-+<U+2212>]?(?:\d{1,3}(?:,\d{3})+|\d+)(?:\.\d+)?)\] against the climatology of its own label
%   matched: Brier skill score 0.6825 [0.6428, 0.7209] against the climatology of its own label
\newcommand{\NumEncPosBSS}{0.6825}  % raw: 0.6825
\newcommand{\NumEncPosLo}{0.6428}  % raw: 0.6428
\newcommand{\NumEncPosHi}{0.7209}  % raw: 0.7209

% docs/RESULTS.md:25, row (2026-10-31 | H1 (encoder arms))
%   pattern: `cheat-stamp` ([-+<U+2212>]?(?:\d{1,3}(?:,\d{3})+|\d+)(?:\.\d+)?) \[([-+<U+2212>]?(?:\d{1,3}(?:,\d{3})+|\d+)(?:\.\d+)?), ([-+<U+2212>]?(?:\d{1,3}(?:,\d{3})+|\d+)(?:\.\d+)?)\] \(interval NOT above
%   matched: `cheat-stamp` 0.0064 [0.0029, 0.0100] (interval NOT above
\newcommand{\NumEncStampBSS}{0.0064}  % raw: 0.0064
\newcommand{\NumEncStampLo}{0.0029}  % raw: 0.0029
\newcommand{\NumEncStampHi}{0.0100}  % raw: 0.0100

% docs/RESULTS.md:25, row (2026-10-31 | H1 (encoder arms))
%   pattern: and below ([-+<U+2212>]?(?:\d{1,3}(?:,\d{3})+|\d+)(?:\.\d+)?): NOT met
%   matched: and below 0.1: NOT met
\newcommand{\NumEncStampFloor}{0.1}  % raw: 0.1

% docs/RESULTS.md:15, row (2026-10-29 | H2)
%   pattern: pooled memory over (\d+) members
%   matched: pooled memory over 233 members
\newcommand{\NumPanelMembers}{233}  % raw: 233

% docs/RESULTS.md:15, row (2026-10-29 | H2)
%   pattern: \((\d{1,3}(?:,\d{3})+|\d+) scored decision hours, (\d{1,3}(?:,\d{3})+|\d+) scored rows, ([-+<U+2212>]?(?:\d{1,3}(?:,\d{3})+|\d+)(?:\.\d+)?) % of the hours\)
%   matched: (50,626 scored decision hours, 1,493,353 scored rows, 96.2 % of the hours)
\newcommand{\NumScoredHours}{50,626}  % raw: 50,626
\newcommand{\NumScoredRows}{1,493,353}  % raw: 1,493,353
\newcommand{\NumScoredShare}{96.2}  % raw: 96.2

% docs/RESULTS.md:15, row (2026-10-29 | H2)
%   pattern: \| (\d+) \(10 before this iteration, plus the 8 arms
%   matched: | 18 (10 before this iteration, plus the 8 arms
\newcommand{\NumRepTwoVariants}{18}  % raw: 18

% docs/RESULTS.md:15, row (2026-10-29 | H2)
%   pattern: BSS of matrix-select ([-+<U+2212>]?(?:\d{1,3}(?:,\d{3})+|\d+)(?:\.\d+)?) \[([-+<U+2212>]?(?:\d{1,3}(?:,\d{3})+|\d+)(?:\.\d+)?), ([-+<U+2212>]?(?:\d{1,3}(?:,\d{3})+|\d+)(?:\.\d+)?)\]
%   matched: BSS of matrix-select 0.0130 [0.0108, 0.0156]
\newcommand{\NumHtwoBSS}{0.0130}  % raw: 0.0130
\newcommand{\NumHtwoBSSlo}{0.0108}  % raw: 0.0108
\newcommand{\NumHtwoBSShi}{0.0156}  % raw: 0.0156

% docs/RESULTS.md:15, row (2026-10-29 | H2)
%   pattern: (\d+) of the (\d+) test months
%   matched: 18 of the 72 test months
\newcommand{\NumPreviewMonths}{18}  % raw: 18
\newcommand{\NumTestMonths}{72}  % raw: 72

% docs/RESULTS.md:15, row (2026-10-29 | H2)
%   pattern: previewed months, apart \(description\): ([-+<U+2212>]?(?:\d{1,3}(?:,\d{3})+|\d+)(?:\.\d+)?) \[([-+<U+2212>]?(?:\d{1,3}(?:,\d{3})+|\d+)(?:\.\d+)?), ([-+<U+2212>]?(?:\d{1,3}(?:,\d{3})+|\d+)(?:\.\d+)?)\] over (\d{1,3}(?:,\d{3})+|\d+) hours
%   matched: previewed months, apart (description): 0.0201 [0.0143, 0.0265] over 12,130 hours
\newcommand{\NumHtwoPrevBSS}{0.0201}  % raw: 0.0201
\newcommand{\NumHtwoPrevLo}{0.0143}  % raw: 0.0143
\newcommand{\NumHtwoPrevHi}{0.0265}  % raw: 0.0265
\newcommand{\NumHtwoPrevHours}{12,130}  % raw: 12,130

% docs/RESULTS.md:15, row (2026-10-29 | H2)
%   pattern: the other (\d+) months: ([-+<U+2212>]?(?:\d{1,3}(?:,\d{3})+|\d+)(?:\.\d+)?) \[([-+<U+2212>]?(?:\d{1,3}(?:,\d{3})+|\d+)(?:\.\d+)?), ([-+<U+2212>]?(?:\d{1,3}(?:,\d{3})+|\d+)(?:\.\d+)?)\] over (\d{1,3}(?:,\d{3})+|\d+) hours
%   matched: the other 54 months: 0.0108 [0.0083, 0.0132] over 38,496 hours
\newcommand{\NumRestMonths}{54}  % raw: 54
\newcommand{\NumHtwoRestBSS}{0.0108}  % raw: 0.0108
\newcommand{\NumHtwoRestLo}{0.0083}  % raw: 0.0083
\newcommand{\NumHtwoRestHi}{0.0132}  % raw: 0.0132
\newcommand{\NumHtwoRestHours}{38,496}  % raw: 38,496

% docs/RESULTS.md:16, row (2026-10-29 | H3)
%   pattern: minus softmax-1y: ([-+<U+2212>]?(?:\d{1,3}(?:,\d{3})+|\d+)(?:\.\d+)?) \(([-+<U+2212>]?(?:\d{1,3}(?:,\d{3})+|\d+)(?:\.\d+)?) % of climatology's log-loss\), DM p ([-+<U+2212>]?(?:\d{1,3}(?:,\d{3})+|\d+)(?:\.\d+)?), bootstrap \[([-+<U+2212>]?(?:\d{1,3}(?:,\d{3})+|\d+)(?:\.\d+)?), ([-+<U+2212>]?(?:\d{1,3}(?:,\d{3})+|\d+)(?:\.\d+)?)\]
%   matched: minus softmax-1y: -0.00891 (-0.8288 % of climatology's log-loss), DM p 0.0000, bootstrap [-0.01101, -0.00691]
\newcommand{\NumHthreeSmDiff}{\textminus{}0.00891}  % raw: -0.00891
\newcommand{\NumHthreeSmPct}{\textminus{}0.8288}  % raw: -0.8288
\newcommand{\NumHthreeSmP}{0.0000}  % raw: 0.0000
\newcommand{\NumHthreeSmLo}{\textminus{}0.01101}  % raw: -0.01101
\newcommand{\NumHthreeSmHi}{\textminus{}0.00691}  % raw: -0.00691

% docs/RESULTS.md:16, row (2026-10-29 | H3)
%   pattern: minus knn-1y: ([-+<U+2212>]?(?:\d{1,3}(?:,\d{3})+|\d+)(?:\.\d+)?) \(([-+<U+2212>]?(?:\d{1,3}(?:,\d{3})+|\d+)(?:\.\d+)?) %\), DM p ([-+<U+2212>]?(?:\d{1,3}(?:,\d{3})+|\d+)(?:\.\d+)?), bootstrap \[([-+<U+2212>]?(?:\d{1,3}(?:,\d{3})+|\d+)(?:\.\d+)?), ([-+<U+2212>]?(?:\d{1,3}(?:,\d{3})+|\d+)(?:\.\d+)?)\]
%   matched: minus knn-1y: -0.00101 (-0.0942 %), DM p 0.2868, bootstrap [-0.00278, +0.00081]
\newcommand{\NumHthreeKnDiff}{\textminus{}0.00101}  % raw: -0.00101
\newcommand{\NumHthreeKnPct}{\textminus{}0.0942}  % raw: -0.0942
\newcommand{\NumHthreeKnP}{0.2868}  % raw: 0.2868
\newcommand{\NumHthreeKnLo}{\textminus{}0.00278}  % raw: -0.00278
\newcommand{\NumHthreeKnHi}{+0.00081}  % raw: +0.00081

% docs/RESULTS.md:17, row (2026-10-29 | H10)
%   pattern: Difference ([-+<U+2212>]?(?:\d{1,3}(?:,\d{3})+|\d+)(?:\.\d+)?) \(([-+<U+2212>]?(?:\d{1,3}(?:,\d{3})+|\d+)(?:\.\d+)?) % of climatology's log-loss;
%   matched: Difference +0.00129 (+0.1199 % of climatology's log-loss;
\newcommand{\NumHtenDiff}{+0.00129}  % raw: +0.00129
\newcommand{\NumHtenPct}{+0.1199}  % raw: +0.1199

% docs/RESULTS.md:17, row (2026-10-29 | H10)
%   pattern: DM p ([-+<U+2212>]?(?:\d{1,3}(?:,\d{3})+|\d+)(?:\.\d+)?), bootstrap \[([-+<U+2212>]?(?:\d{1,3}(?:,\d{3})+|\d+)(?:\.\d+)?), ([-+<U+2212>]?(?:\d{1,3}(?:,\d{3})+|\d+)(?:\.\d+)?)\], does not pass\. Identical hours ([-+<U+2212>]?(?:\d{1,3}(?:,\d{3})+|\d+)(?:\.\d+)?) %
%   matched: DM p 0.0060, bootstrap [+0.00049, +0.00228], does not pass. Identical hours 31.9 %
\newcommand{\NumHtenP}{0.0060}  % raw: 0.0060
\newcommand{\NumHtenLo}{+0.00049}  % raw: +0.00049
\newcommand{\NumHtenHi}{+0.00228}  % raw: +0.00228
\newcommand{\NumHtenIdent}{31.9}  % raw: 31.9

% docs/RESULTS.md:18, row (2026-10-29 | H11)
%   pattern: Difference ([-+<U+2212>]?(?:\d{1,3}(?:,\d{3})+|\d+)(?:\.\d+)?) \(([-+<U+2212>]?(?:\d{1,3}(?:,\d{3})+|\d+)(?:\.\d+)?) % of climatology's log-loss\), DM p ([-+<U+2212>]?(?:\d{1,3}(?:,\d{3})+|\d+)(?:\.\d+)?), bootstrap \[([-+<U+2212>]?(?:\d{1,3}(?:,\d{3})+|\d+)(?:\.\d+)?), ([-+<U+2212>]?(?:\d{1,3}(?:,\d{3})+|\d+)(?:\.\d+)?)\]
%   matched: Difference -0.00060 (-0.0555 % of climatology's log-loss), DM p 0.0725, bootstrap [-0.00122, +0.00001]
\newcommand{\NumHelevenDiff}{\textminus{}0.00060}  % raw: -0.00060
\newcommand{\NumHelevenPct}{\textminus{}0.0555}  % raw: -0.0555
\newcommand{\NumHelevenP}{0.0725}  % raw: 0.0725
\newcommand{\NumHelevenLo}{\textminus{}0.00122}  % raw: -0.00122
\newcommand{\NumHelevenHi}{+0.00001}  % raw: +0.00001

% docs/RESULTS.md:19, row (2026-10-29 | H12)
%   pattern: Difference ([-+<U+2212>]?(?:\d{1,3}(?:,\d{3})+|\d+)(?:\.\d+)?) \(([-+<U+2212>]?(?:\d{1,3}(?:,\d{3})+|\d+)(?:\.\d+)?) % of climatology's log-loss\), DM p ([-+<U+2212>]?(?:\d{1,3}(?:,\d{3})+|\d+)(?:\.\d+)?), bootstrap \[([-+<U+2212>]?(?:\d{1,3}(?:,\d{3})+|\d+)(?:\.\d+)?), ([-+<U+2212>]?(?:\d{1,3}(?:,\d{3})+|\d+)(?:\.\d+)?)\]
%   matched: Difference -0.00002 (-0.0022 % of climatology's log-loss), DM p 0.9440, bootstrap [-0.00071, +0.00062]
\newcommand{\NumHtwelveDiff}{\textminus{}0.00002}  % raw: -0.00002
\newcommand{\NumHtwelvePct}{\textminus{}0.0022}  % raw: -0.0022
\newcommand{\NumHtwelveP}{0.9440}  % raw: 0.9440
\newcommand{\NumHtwelveLo}{\textminus{}0.00071}  % raw: -0.00071
\newcommand{\NumHtwelveHi}{+0.00062}  % raw: +0.00062

% docs/RESULTS.md:19, row (2026-10-29 | H12)
%   pattern: margin ([-+<U+2212>]?(?:\d{1,3}(?:,\d{3})+|\d+)(?:\.\d+)?) x climatology's log-loss = ([-+<U+2212>]?(?:\d{1,3}(?:,\d{3})+|\d+)(?:\.\d+)?)
%   matched: margin 0.005 x climatology's log-loss = 0.00538
\newcommand{\NumMarginFactor}{0.005}  % raw: 0.005
\newcommand{\NumMargin}{0.00538}  % raw: 0.00538

% docs/RESULTS.md:19, row (2026-10-29 | H12)
%   pattern: one-sided 95 % upper bound ([-+<U+2212>]?(?:\d{1,3}(?:,\d{3})+|\d+)(?:\.\d+)?); bootstrap 95th percentile ([-+<U+2212>]?(?:\d{1,3}(?:,\d{3})+|\d+)(?:\.\d+)?); both below
%   matched: one-sided 95 % upper bound +0.00053; bootstrap 95th percentile +0.00052; both below
\newcommand{\NumHtwelveUB}{+0.00053}  % raw: +0.00053
\newcommand{\NumHtwelveQ}{+0.00052}  % raw: +0.00052

% docs/RESULTS.md:19, row (2026-10-29 | H12)
%   pattern: Fits: (\d+), against the limit of (\d+) \((\d+) alarms over (\d{1,3}(?:,\d{3})+|\d+) daily values\)
%   matched: Fits: 13, against the limit of 36 (12 alarms over 2,108 daily values)
\newcommand{\NumHtwelveFits}{13}  % raw: 13
\newcommand{\NumHtwelveLimit}{36}  % raw: 36
\newcommand{\NumHtwelveAlarms}{12}  % raw: 12
\newcommand{\NumHtwelveDays}{2,108}  % raw: 2,108

% docs/RESULTS.md:19, row (2026-10-29 | H12)
%   pattern: Page CUSUM alarm, k = ([-+<U+2212>]?(?:\d{1,3}(?:,\d{3})+|\d+)(?:\.\d+)?), threshold ([-+<U+2212>]?(?:\d{1,3}(?:,\d{3})+|\d+)(?:\.\d+)?), ARL (\d+) days
%   matched: Page CUSUM alarm, k = 0.5, threshold 4.761, ARL 365 days
\newcommand{\NumCusumK}{0.5}  % raw: 0.5
\newcommand{\NumCusumH}{4.761}  % raw: 4.761
\newcommand{\NumCusumARL}{365}  % raw: 365

% docs/RESULTS.md:20, row (2026-10-29 | Report n<U+00B0>2 (descriptive))
%   pattern: (?<=\s)climatology-panel ([-+<U+2212>]?(?:\d{1,3}(?:,\d{3})+|\d+)(?:\.\d+)?) \[([-+<U+2212>]?(?:\d{1,3}(?:,\d{3})+|\d+)(?:\.\d+)?), ([-+<U+2212>]?(?:\d{1,3}(?:,\d{3})+|\d+)(?:\.\d+)?)\]
%   matched: climatology-panel 1.0756 [1.0726, 1.0786]
\newcommand{\NumLLclim}{1.0756}  % raw: 1.0756
\newcommand{\NumLLclimLo}{1.0726}  % raw: 1.0726
\newcommand{\NumLLclimHi}{1.0786}  % raw: 1.0786

% docs/RESULTS.md:20, row (2026-10-29 | Report n<U+00B0>2 (descriptive))
%   pattern: (?<=\s)matrix-select ([-+<U+2212>]?(?:\d{1,3}(?:,\d{3})+|\d+)(?:\.\d+)?) \[([-+<U+2212>]?(?:\d{1,3}(?:,\d{3})+|\d+)(?:\.\d+)?), ([-+<U+2212>]?(?:\d{1,3}(?:,\d{3})+|\d+)(?:\.\d+)?)\]
%   matched: matrix-select 1.0650 [1.0611, 1.0690]
\newcommand{\NumLLsel}{1.0650}  % raw: 1.0650
\newcommand{\NumLLselLo}{1.0611}  % raw: 1.0611
\newcommand{\NumLLselHi}{1.0690}  % raw: 1.0690

% docs/RESULTS.md:20, row (2026-10-29 | Report n<U+00B0>2 (descriptive))
%   pattern: (?<=\s)matrix-1y ([-+<U+2212>]?(?:\d{1,3}(?:,\d{3})+|\d+)(?:\.\d+)?) \[([-+<U+2212>]?(?:\d{1,3}(?:,\d{3})+|\d+)(?:\.\d+)?), ([-+<U+2212>]?(?:\d{1,3}(?:,\d{3})+|\d+)(?:\.\d+)?)\]
%   matched: matrix-1y 1.0651 [1.0611, 1.0697]
\newcommand{\NumLLoney}{1.0651}  % raw: 1.0651
\newcommand{\NumLLoneyLo}{1.0611}  % raw: 1.0611
\newcommand{\NumLLoneyHi}{1.0697}  % raw: 1.0697

% docs/RESULTS.md:20, row (2026-10-29 | Report n<U+00B0>2 (descriptive))
%   pattern: (?<=\s)softmax-1y ([-+<U+2212>]?(?:\d{1,3}(?:,\d{3})+|\d+)(?:\.\d+)?) \[([-+<U+2212>]?(?:\d{1,3}(?:,\d{3})+|\d+)(?:\.\d+)?), ([-+<U+2212>]?(?:\d{1,3}(?:,\d{3})+|\d+)(?:\.\d+)?)\]
%   matched: softmax-1y 1.0740 [1.0710, 1.0771]
\newcommand{\NumLLsoft}{1.0740}  % raw: 1.0740
\newcommand{\NumLLsoftLo}{1.0710}  % raw: 1.0710
\newcommand{\NumLLsoftHi}{1.0771}  % raw: 1.0771

% docs/RESULTS.md:20, row (2026-10-29 | Report n<U+00B0>2 (descriptive))
%   pattern: (?<=\s)knn-1y ([-+<U+2212>]?(?:\d{1,3}(?:,\d{3})+|\d+)(?:\.\d+)?) \[([-+<U+2212>]?(?:\d{1,3}(?:,\d{3})+|\d+)(?:\.\d+)?), ([-+<U+2212>]?(?:\d{1,3}(?:,\d{3})+|\d+)(?:\.\d+)?)\]
%   matched: knn-1y 1.0661 [1.0620, 1.0701]
\newcommand{\NumLLknn}{1.0661}  % raw: 1.0661
\newcommand{\NumLLknnLo}{1.0620}  % raw: 1.0620
\newcommand{\NumLLknnHi}{1.0701}  % raw: 1.0701

% docs/RESULTS.md:20, row (2026-10-29 | Report n<U+00B0>2 (descriptive))
%   pattern: (?<=\s)matrix-adaptive ([-+<U+2212>]?(?:\d{1,3}(?:,\d{3})+|\d+)(?:\.\d+)?) \[([-+<U+2212>]?(?:\d{1,3}(?:,\d{3})+|\d+)(?:\.\d+)?), ([-+<U+2212>]?(?:\d{1,3}(?:,\d{3})+|\d+)(?:\.\d+)?)\]
%   matched: matrix-adaptive 1.0663 [1.0626, 1.0702]
\newcommand{\NumLLadapt}{1.0663}  % raw: 1.0663
\newcommand{\NumLLadaptLo}{1.0626}  % raw: 1.0626
\newcommand{\NumLLadaptHi}{1.0702}  % raw: 1.0702

% docs/RESULTS.md:20, row (2026-10-29 | Report n<U+00B0>2 (descriptive))
%   pattern: (?<=\s)matrix-stack ([-+<U+2212>]?(?:\d{1,3}(?:,\d{3})+|\d+)(?:\.\d+)?) \[([-+<U+2212>]?(?:\d{1,3}(?:,\d{3})+|\d+)(?:\.\d+)?), ([-+<U+2212>]?(?:\d{1,3}(?:,\d{3})+|\d+)(?:\.\d+)?)\]
%   matched: matrix-stack 1.0644 [1.0603, 1.0686]
\newcommand{\NumLLstack}{1.0644}  % raw: 1.0644
\newcommand{\NumLLstackLo}{1.0603}  % raw: 1.0603
\newcommand{\NumLLstackHi}{1.0686}  % raw: 1.0686

% docs/RESULTS.md:20, row (2026-10-29 | Report n<U+00B0>2 (descriptive))
%   pattern: (?<=\s)matrix-cusum ([-+<U+2212>]?(?:\d{1,3}(?:,\d{3})+|\d+)(?:\.\d+)?) \[([-+<U+2212>]?(?:\d{1,3}(?:,\d{3})+|\d+)(?:\.\d+)?), ([-+<U+2212>]?(?:\d{1,3}(?:,\d{3})+|\d+)(?:\.\d+)?)\]
%   matched: matrix-cusum 1.0650 [1.0609, 1.0690]
\newcommand{\NumLLcusum}{1.0650}  % raw: 1.0650
\newcommand{\NumLLcusumLo}{1.0609}  % raw: 1.0609
\newcommand{\NumLLcusumHi}{1.0690}  % raw: 1.0690

% docs/RESULTS.md:20, row (2026-10-29 | Report n<U+00B0>2 (descriptive))
%   pattern: BSS against climatology: ([-+<U+2212>]?(?:\d{1,3}(?:,\d{3})+|\d+)(?:\.\d+)?), ([-+<U+2212>]?(?:\d{1,3}(?:,\d{3})+|\d+)(?:\.\d+)?), ([-+<U+2212>]?(?:\d{1,3}(?:,\d{3})+|\d+)(?:\.\d+)?), ([-+<U+2212>]?(?:\d{1,3}(?:,\d{3})+|\d+)(?:\.\d+)?), ([-+<U+2212>]?(?:\d{1,3}(?:,\d{3})+|\d+)(?:\.\d+)?), ([-+<U+2212>]?(?:\d{1,3}(?:,\d{3})+|\d+)(?:\.\d+)?) and ([-+<U+2212>]?(?:\d{1,3}(?:,\d{3})+|\d+)(?:\.\d+)?) for the seven other arms
%   matched: BSS against climatology: 0.0130, 0.0134, 0.0032, 0.0109, 0.0115, 0.0131 and 0.0130 for the seven other arms
\newcommand{\NumBSSsel}{0.0130}  % raw: 0.0130
\newcommand{\NumBSSoney}{0.0134}  % raw: 0.0134
\newcommand{\NumBSSsoft}{0.0032}  % raw: 0.0032
\newcommand{\NumBSSknn}{0.0109}  % raw: 0.0109
\newcommand{\NumBSSadapt}{0.0115}  % raw: 0.0115
\newcommand{\NumBSSstack}{0.0131}  % raw: 0.0131
\newcommand{\NumBSScusum}{0.0130}  % raw: 0.0130

% docs/RESULTS.md:20, row (2026-10-29 | Report n<U+00B0>2 (descriptive))
%   pattern: `--check` reproduces its (\d+) files byte for byte
%   matched: `--check` reproduces its 3 files byte for byte
\newcommand{\NumRepTwoCheckFiles}{3}  % raw: 3

% docs/RESULTS.md:21, row (2026-10-30 | Report n<U+00B0>2 (re-run, descriptive))
%   pattern: `--check` reproduces its (\d+) files byte for byte
%   matched: `--check` reproduces its 4 files byte for byte
\newcommand{\NumRepTwoRerunFiles}{4}  % raw: 4

% docs/RESULTS.md:23, row (2026-10-30 | Report n<U+00B0>2 (WSL2 re-run, descriptive))
%   pattern: Python ([\d.]+) \(Windows: ([\d.]+)\)
%   matched: Python 3.12.14 (Windows: 3.12.10)
\newcommand{\NumPyWSL}{3.12.14}  % raw: 3.12.14
\newcommand{\NumPyWin}{3.12.10}  % raw: 3.12.10

% docs/RESULTS.md:23, row (2026-10-30 | Report n<U+00B0>2 (WSL2 re-run, descriptive))
%   pattern: Probabilities differ in (\d+) of (\d+) arms, by at most (\d\.\d <U+00D7> 10[<U+207B><U+00B9><U+00B2><U+00B3><U+2074><U+2075><U+2076><U+2077><U+2078><U+2079><U+2070>]+)
%   matched: Probabilities differ in 6 of 8 arms, by at most 1.6 <U+00D7> 10<U+207B><U+00B9><U+2075>
\newcommand{\NumPlatProbArms}{6}  % raw: 6
\newcommand{\NumPlatArms}{8}  % raw: 8
\newcommand{\NumPlatProbMax}{\ensuremath{1.6\times 10^{-15}}}  % raw: 1.6 <U+00D7> 10<U+207B><U+00B9><U+2075>

% docs/RESULTS.md:23, row (2026-10-30 | Report n<U+00B0>2 (WSL2 re-run, descriptive))
%   pattern: (\d+) of (\d{1,3}(?:,\d{3})+|\d+) numbers of `report2.json` differ, by at most (\d\.\d <U+00D7> 10[<U+207B><U+00B9><U+00B2><U+00B3><U+2074><U+2075><U+2076><U+2077><U+2078><U+2079><U+2070>]+) relative
%   matched: 162 of 2,249 numbers of `report2.json` differ, by at most 5.2 <U+00D7> 10<U+207B><U+00B9><U+2074> relative
\newcommand{\NumPlatNums}{162}  % raw: 162
\newcommand{\NumPlatNumsOf}{2,249}  % raw: 2,249
\newcommand{\NumPlatRelMax}{\ensuremath{5.2\times 10^{-14}}}  % raw: 5.2 <U+00D7> 10<U+207B><U+00B9><U+2074>

% docs/RESULTS.md:24, row (2026-10-30 | Size and power of the pass rule (descriptive))
%   pattern: L <U+2208> \{(\d+), (\d+), (\d+), (\d+)\} hours
%   matched: L <U+2208> {24, 48, 168, 720} hours
\newcommand{\NumPowLa}{24}  % raw: 24
\newcommand{\NumPowLb}{48}  % raw: 48
\newcommand{\NumPowLc}{168}  % raw: 168
\newcommand{\NumPowLd}{720}  % raw: 720

% docs/RESULTS.md:24, row (2026-10-30 | Size and power of the pass rule (descriptive))
%   pattern: (\d{1,3}(?:,\d{3})+|\d+) replicates per L for size
%   matched: 2,000 replicates per L for size
\newcommand{\NumPowSizeReps}{2,000}  % raw: 2,000

% docs/RESULTS.md:24, row (2026-10-30 | Size and power of the pass rule (descriptive))
%   pattern: grid of (\d+) points with (\d{1,3}(?:,\d{3})+|\d+) replicates each
%   matched: grid of 10 points with 1,000 replicates each
\newcommand{\NumPowGrid}{10}  % raw: 10
\newcommand{\NumPowGridReps}{1,000}  % raw: 1,000

% docs/RESULTS.md:24, row (2026-10-30 | Size and power of the pass rule (descriptive))
%   pattern: concentrated in (\d+) % of (\d+)-hour blocks \(L = (\d+)\)
%   matched: concentrated in 10 % of 720-hour blocks (L = 168)
\newcommand{\NumPowConcPct}{10}  % raw: 10
\newcommand{\NumPowConcBlock}{720}  % raw: 720
\newcommand{\NumPowConcL}{168}  % raw: 168

% docs/RESULTS.md:24, row (2026-10-30 | Size and power of the pass rule (descriptive))
%   pattern: flip seed (\d+); (\d{1,3}(?:,\d{3})+|\d+) replicates\.
%   matched: flip seed 20261030; 244,000 replicates.
\newcommand{\NumPowSeed}{20261030}  % raw: 20261030
\newcommand{\NumPowReps}{244,000}  % raw: 244,000

% docs/RESULTS.md:24, row (2026-10-30 | Size and power of the pass rule (descriptive))
%   pattern: two-sided size is ([-+<U+2212>]?(?:\d{1,3}(?:,\d{3})+|\d+)(?:\.\d+)?) % to ([-+<U+2212>]?(?:\d{1,3}(?:,\d{3})+|\d+)(?:\.\d+)?) % for every comparison and L: none above ([-+<U+2212>]?(?:\d{1,3}(?:,\d{3})+|\d+)(?:\.\d+)?) %
%   matched: two-sided size is 2.2 % to 6.4 % for every comparison and L: none above 7.5 %
\newcommand{\NumPowSizeLo}{2.2}  % raw: 2.2
\newcommand{\NumPowSizeHi}{6.4}  % raw: 6.4
\newcommand{\NumPowSizeCap}{7.5}  % raw: 7.5

% docs/RESULTS.md:24, row (2026-10-30 | Size and power of the pass rule (descriptive))
%   pattern: registered direction, ([-+<U+2212>]?(?:\d{1,3}(?:,\d{3})+|\d+)(?:\.\d+)?) % to ([-+<U+2212>]?(?:\d{1,3}(?:,\d{3})+|\d+)(?:\.\d+)?) %
%   matched: registered direction, 1.1 % to 2.8 %
\newcommand{\NumPowDirLo}{1.1}  % raw: 1.1
\newcommand{\NumPowDirHi}{2.8}  % raw: 2.8

% docs/RESULTS.md:24, row (2026-10-30 | Size and power of the pass rule (descriptive))
%   pattern: bootstrap alone reaches ([-+<U+2212>]?(?:\d{1,3}(?:,\d{3})+|\d+)(?:\.\d+)?) % \(H10, L = (\d+) and (\d+)\)
%   matched: bootstrap alone reaches 7.8 % (H10, L = 168 and 720)
\newcommand{\NumPowBootMax}{7.8}  % raw: 7.8
\newcommand{\NumPowBootLa}{168}  % raw: 168
\newcommand{\NumPowBootLb}{720}  % raw: 720

% docs/RESULTS.md:24, row (2026-10-30 | Size and power of the pass rule (descriptive))
%   pattern: H12 holds at its margin ([-+<U+2212>]?(?:\d{1,3}(?:,\d{3})+|\d+)(?:\.\d+)?) % to ([-+<U+2212>]?(?:\d{1,3}(?:,\d{3})+|\d+)(?:\.\d+)?) %
%   matched: H12 holds at its margin 4.4 % to 6.4 %
\newcommand{\NumPowTwelveLo}{4.4}  % raw: 4.4
\newcommand{\NumPowTwelveHi}{6.4}  % raw: 6.4

% docs/RESULTS.md:24, row (2026-10-30 | Size and power of the pass rule (descriptive))
%   pattern: smallest <U+03B4> with (\d+) % power
%   matched: smallest <U+03B4> with 80 % power
\newcommand{\NumPowTarget}{80}  % raw: 80

% docs/RESULTS.md:24, row (2026-10-30 | Size and power of the pass rule (descriptive))
%   pattern: H3 vs softmax ([-+<U+2212>]?(?:\d{1,3}(?:,\d{3})+|\d+)(?:\.\d+)?) % to ([-+<U+2212>]?(?:\d{1,3}(?:,\d{3})+|\d+)(?:\.\d+)?) % \(([-+<U+2212>]?(?:\d{1,3}(?:,\d{3})+|\d+)(?:\.\d+)?) %\)
%   matched: H3 vs softmax 0.26 % to 0.31 % (0.35 %)
\newcommand{\NumPowSmLo}{0.26}  % raw: 0.26
\newcommand{\NumPowSmHi}{0.31}  % raw: 0.31
\newcommand{\NumPowSmConc}{0.35}  % raw: 0.35

% docs/RESULTS.md:24, row (2026-10-30 | Size and power of the pass rule (descriptive))
%   pattern: H3 vs kNN ([-+<U+2212>]?(?:\d{1,3}(?:,\d{3})+|\d+)(?:\.\d+)?) % to ([-+<U+2212>]?(?:\d{1,3}(?:,\d{3})+|\d+)(?:\.\d+)?) % \(([-+<U+2212>]?(?:\d{1,3}(?:,\d{3})+|\d+)(?:\.\d+)?) %\)
%   matched: H3 vs kNN 0.23 % to 0.25 % (0.26 %)
\newcommand{\NumPowKnLo}{0.23}  % raw: 0.23
\newcommand{\NumPowKnHi}{0.25}  % raw: 0.25
\newcommand{\NumPowKnConc}{0.26}  % raw: 0.26

% docs/RESULTS.md:24, row (2026-10-30 | Size and power of the pass rule (descriptive))
%   pattern: H10 ([-+<U+2212>]?(?:\d{1,3}(?:,\d{3})+|\d+)(?:\.\d+)?) % to ([-+<U+2212>]?(?:\d{1,3}(?:,\d{3})+|\d+)(?:\.\d+)?) % \(([-+<U+2212>]?(?:\d{1,3}(?:,\d{3})+|\d+)(?:\.\d+)?) %\)
%   matched: H10 0.12 % to 0.14 % (0.15 %)
\newcommand{\NumPowTenLo}{0.12}  % raw: 0.12
\newcommand{\NumPowTenHi}{0.14}  % raw: 0.14
\newcommand{\NumPowTenConc}{0.15}  % raw: 0.15

% docs/RESULTS.md:24, row (2026-10-30 | Size and power of the pass rule (descriptive))
%   pattern: H11 ([-+<U+2212>]?(?:\d{1,3}(?:,\d{3})+|\d+)(?:\.\d+)?) % to ([-+<U+2212>]?(?:\d{1,3}(?:,\d{3})+|\d+)(?:\.\d+)?) % \(([-+<U+2212>]?(?:\d{1,3}(?:,\d{3})+|\d+)(?:\.\d+)?) %\)
%   matched: H11 0.08 % to 0.09 % (0.10 %)
\newcommand{\NumPowElLo}{0.08}  % raw: 0.08
\newcommand{\NumPowElHi}{0.09}  % raw: 0.09
\newcommand{\NumPowElConc}{0.10}  % raw: 0.10

% docs/RESULTS.md:24, row (2026-10-30 | Size and power of the pass rule (descriptive))
%   pattern: `knn-1y`, ([-+<U+2212>]?(?:\d{1,3}(?:,\d{3})+|\d+)(?:\.\d+)?) %, has (\d+) % to (\d+) % power
%   matched: `knn-1y`, <U+2212>0.094 %, has 19 % to 24 % power
\newcommand{\NumPowKnObs}{\textminus{}0.094}  % raw: <U+2212>0.094
\newcommand{\NumPowKnObsLo}{19}  % raw: 19
\newcommand{\NumPowKnObsHi}{24}  % raw: 24

% docs/RESULTS.md:24, row (2026-10-30 | Size and power of the pass rule (descriptive))
%   pattern: H11's ([-+<U+2212>]?(?:\d{1,3}(?:,\d{3})+|\d+)(?:\.\d+)?) %, (\d+) % to (\d+) %
%   matched: H11's <U+2212>0.056 %, 42 % to 51 %
\newcommand{\NumPowElObs}{\textminus{}0.056}  % raw: <U+2212>0.056
\newcommand{\NumPowElObsLo}{42}  % raw: 42
\newcommand{\NumPowElObsHi}{51}  % raw: 51

% docs/RESULTS.md:24, row (2026-10-30 | Size and power of the pass rule (descriptive))
%   pattern: first run's facts: (\d{1,3}(?:,\d{3})+|\d+) s, (\d+) workers
%   matched: first run's facts: 5,198 s, 8 workers
\newcommand{\NumPowSecs}{5,198}  % raw: 5,198
\newcommand{\NumPowWorkers}{8}  % raw: 8

% docs/RESULTS.md:26, row (2026-10-31 | H4)
%   pattern: \| (\d+) \(18 before, plus the 3 encoder arms\)
%   matched: | 21 (18 before, plus the 3 encoder arms)
\newcommand{\NumRepThreeVariants}{21}  % raw: 21

% docs/RESULTS.md:26, row (2026-10-31 | H4)
%   pattern: (\d+) seeds per arm, (\d+) yearly fits per seed
%   matched: 3 seeds per arm, 6 yearly fits per seed
\newcommand{\NumSeeds}{3}  % raw: 3
\newcommand{\NumYearlyFits}{6}  % raw: 6

% docs/RESULTS.md:26, row (2026-10-31 | H4)
%   pattern: BSS `enc-ols` ([-+<U+2212>]?(?:\d{1,3}(?:,\d{3})+|\d+)(?:\.\d+)?) \[([-+<U+2212>]?(?:\d{1,3}(?:,\d{3})+|\d+)(?:\.\d+)?), ([-+<U+2212>]?(?:\d{1,3}(?:,\d{3})+|\d+)(?:\.\d+)?)\], `enc-softmax` ([-+<U+2212>]?(?:\d{1,3}(?:,\d{3})+|\d+)(?:\.\d+)?) \[([-+<U+2212>]?(?:\d{1,3}(?:,\d{3})+|\d+)(?:\.\d+)?), ([-+<U+2212>]?(?:\d{1,3}(?:,\d{3})+|\d+)(?:\.\d+)?)\]
%   matched: BSS `enc-ols` 0.0110 [0.0087, 0.0133], `enc-softmax` 0.0079 [0.0059, 0.0098]
\newcommand{\NumHfourBSSols}{0.0110}  % raw: 0.0110
\newcommand{\NumHfourBSSolsLo}{0.0087}  % raw: 0.0087
\newcommand{\NumHfourBSSolsHi}{0.0133}  % raw: 0.0133
\newcommand{\NumHfourBSSsm}{0.0079}  % raw: 0.0079
\newcommand{\NumHfourBSSsmLo}{0.0059}  % raw: 0.0059
\newcommand{\NumHfourBSSsmHi}{0.0098}  % raw: 0.0098

% docs/RESULTS.md:26, row (2026-10-31 | H4)
%   pattern: Difference ([-+<U+2212>]?(?:\d{1,3}(?:,\d{3})+|\d+)(?:\.\d+)?) \(([-+<U+2212>]?(?:\d{1,3}(?:,\d{3})+|\d+)(?:\.\d+)?) % of climatology's log-loss [\d.]+\); margin ([-+<U+2212>]?(?:\d{1,3}(?:,\d{3})+|\d+)(?:\.\d+)?); one-sided 95 % upper bound ([-+<U+2212>]?(?:\d{1,3}(?:,\d{3})+|\d+)(?:\.\d+)?); bootstrap 95th percentile ([-+<U+2212>]?(?:\d{1,3}(?:,\d{3})+|\d+)(?:\.\d+)?);
%   matched: Difference -0.00269 (-0.2498 % of climatology's log-loss 1.0756); margin 0.00538; one-sided 95 % upper bound -0.00182; bootstrap 95th percentile -0.00181;
\newcommand{\NumHfourDiff}{\textminus{}0.00269}  % raw: -0.00269
\newcommand{\NumHfourPct}{\textminus{}0.2498}  % raw: -0.2498
\newcommand{\NumHfourMargin}{0.00538}  % raw: 0.00538
\newcommand{\NumHfourUB}{\textminus{}0.00182}  % raw: -0.00182
\newcommand{\NumHfourQ}{\textminus{}0.00181}  % raw: -0.00181

% docs/RESULTS.md:26, row (2026-10-31 | H4)
%   pattern: Bootstrap 95 % interval \[([-+<U+2212>]?(?:\d{1,3}(?:,\d{3})+|\d+)(?:\.\d+)?), ([-+<U+2212>]?(?:\d{1,3}(?:,\d{3})+|\d+)(?:\.\d+)?)\], DM p ([\d.]+e-\d+)
%   matched: Bootstrap 95 % interval [-0.00369, -0.00161], DM p 3.0e-07
\newcommand{\NumHfourLo}{\textminus{}0.00369}  % raw: -0.00369
\newcommand{\NumHfourHi}{\textminus{}0.00161}  % raw: -0.00161
\newcommand{\NumHfourP}{\ensuremath{3.0\times 10^{-7}}}  % raw: 3.0e-07

% docs/RESULTS.md:26, row (2026-10-31 | H4)
%   pattern: upper bounds ([-+<U+2212>]?(?:\d{1,3}(?:,\d{3})+|\d+)(?:\.\d+)?), ([-+<U+2212>]?(?:\d{1,3}(?:,\d{3})+|\d+)(?:\.\d+)?), ([-+<U+2212>]?(?:\d{1,3}(?:,\d{3})+|\d+)(?:\.\d+)?)\)
%   matched: upper bounds -0.00172, -0.00046, -0.00068)
\newcommand{\NumHfourSeedA}{\textminus{}0.00172}  % raw: -0.00172
\newcommand{\NumHfourSeedB}{\textminus{}0.00046}  % raw: -0.00046
\newcommand{\NumHfourSeedC}{\textminus{}0.00068}  % raw: -0.00068

% docs/RESULTS.md:26, row (2026-10-31 | H4)
%   pattern: has upper bound ([-+<U+2212>]?(?:\d{1,3}(?:,\d{3})+|\d+)(?:\.\d+)?), also below the margin
%   matched: has upper bound +0.00355, also below the margin
\newcommand{\NumHfourRev}{+0.00355}  % raw: +0.00355

% docs/RESULTS.md:27, row (2026-10-31 | H5)
%   pattern: Difference ([-+<U+2212>]?(?:\d{1,3}(?:,\d{3})+|\d+)(?:\.\d+)?) \(([-+<U+2212>]?(?:\d{1,3}(?:,\d{3})+|\d+)(?:\.\d+)?) % of climatology's log-loss; positive: `enc-ols` worse\), DM p ([-+<U+2212>]?(?:\d{1,3}(?:,\d{3})+|\d+)(?:\.\d+)?), bootstrap 95 % interval \[([-+<U+2212>]?(?:\d{1,3}(?:,\d{3})+|\d+)(?:\.\d+)?), ([-+<U+2212>]?(?:\d{1,3}(?:,\d{3})+|\d+)(?:\.\d+)?)\]
%   matched: Difference +0.00082 (+0.0759 % of climatology's log-loss; positive: `enc-ols` worse), DM p 0.5198, bootstrap 95 % interval [-0.00147, +0.00313]
\newcommand{\NumHfiveDiff}{+0.00082}  % raw: +0.00082
\newcommand{\NumHfivePct}{+0.0759}  % raw: +0.0759
\newcommand{\NumHfiveP}{0.5198}  % raw: 0.5198
\newcommand{\NumHfiveLo}{\textminus{}0.00147}  % raw: -0.00147
\newcommand{\NumHfiveHi}{+0.00313}  % raw: +0.00313

% docs/RESULTS.md:27, row (2026-10-31 | H5)
%   pattern: Log-loss ([-+<U+2212>]?(?:\d{1,3}(?:,\d{3})+|\d+)(?:\.\d+)?) for `enc-ols` against ([-+<U+2212>]?(?:\d{1,3}(?:,\d{3})+|\d+)(?:\.\d+)?) for `matrix-select`
%   matched: Log-loss 1.0659 for `enc-ols` against 1.0650 for `matrix-select`
\newcommand{\NumHfiveLLols}{1.0659}  % raw: 1.0659
\newcommand{\NumHfiveLLsel}{1.0650}  % raw: 1.0650

% docs/RESULTS.md:28, row (2026-10-31 | H13)
%   pattern: Registered ratio ([-+<U+2212>]?(?:\d{1,3}(?:,\d{3})+|\d+)(?:\.\d+)?) \(paired-median ratio [\d.]+\); bound ([-+<U+2212>]?(?:\d{1,3}(?:,\d{3})+|\d+)(?:\.\d+)?)
%   matched: Registered ratio 0.4031 (paired-median ratio 0.4038); bound 0.9
\newcommand{\NumHthirteenRatio}{0.4031}  % raw: 0.4031
\newcommand{\NumHthirteenBound}{0.9}  % raw: 0.9

% docs/RESULTS.md:28, row (2026-10-31 | H13)
%   pattern: : (\d+) repeats cycling the arms
%   matched: : 10 repeats cycling the arms
\newcommand{\NumHthirteenRepeats}{10}  % raw: 10

% docs/RESULTS.md:28, row (2026-10-31 | H13)
%   pattern: (\d{1,3}(?:,\d{3})+|\d+) MiB reserved on a card of (\d{1,3}(?:,\d{3})+|\d+) MiB
%   matched: 12,926 MiB reserved on a card of 12,288 MiB
\newcommand{\NumCacheReserved}{12,926}  % raw: 12,926
\newcommand{\NumCardTotal}{12,288}  % raw: 12,288

% docs/RESULTS.md:28, row (2026-10-31 | H13)
%   pattern: the ratio is ([-+<U+2212>]?(?:\d{1,3}(?:,\d{3})+|\d+)(?:\.\d+)?) as the code stands and ([-+<U+2212>]?(?:\d{1,3}(?:,\d{3})+|\d+)(?:\.\d+)?) with that cache emptied
%   matched: the ratio is 0.3739 as the code stands and 1.0732 with that cache emptied
\newcommand{\NumRatioAsIs}{0.3739}  % raw: 0.3739
\newcommand{\NumRatioEmptied}{1.0732}  % raw: 1.0732

% docs/RESULTS.md:28, row (2026-10-31 | H13)
%   pattern: Synthetic rows give ([-+<U+2212>]?(?:\d{1,3}(?:,\d{3})+|\d+)(?:\.\d+)?) to ([-+<U+2212>]?(?:\d{1,3}(?:,\d{3})+|\d+)(?:\.\d+)?) on that
%   matched: Synthetic rows give 1.057 to 1.064 on that
\newcommand{\NumRatioSynthLo}{1.057}  % raw: 1.057
\newcommand{\NumRatioSynthHi}{1.064}  % raw: 1.064

% docs/RESULTS.md:29, row (2026-10-31 | Report n<U+00B0>3 (descriptive))
%   pattern: `enc-ols` log-loss ([-+<U+2212>]?(?:\d{1,3}(?:,\d{3})+|\d+)(?:\.\d+)?) \[([-+<U+2212>]?(?:\d{1,3}(?:,\d{3})+|\d+)(?:\.\d+)?), ([-+<U+2212>]?(?:\d{1,3}(?:,\d{3})+|\d+)(?:\.\d+)?)\], BSS ([-+<U+2212>]?(?:\d{1,3}(?:,\d{3})+|\d+)(?:\.\d+)?) \[([-+<U+2212>]?(?:\d{1,3}(?:,\d{3})+|\d+)(?:\.\d+)?), ([-+<U+2212>]?(?:\d{1,3}(?:,\d{3})+|\d+)(?:\.\d+)?)\]
%   matched: `enc-ols` log-loss 1.0659 [1.0623, 1.0696], BSS 0.0110 [0.0087, 0.0133]
\newcommand{\NumLLols}{1.0659}  % raw: 1.0659
\newcommand{\NumLLolsLo}{1.0623}  % raw: 1.0623
\newcommand{\NumLLolsHi}{1.0696}  % raw: 1.0696
\newcommand{\NumBSSols}{0.0110}  % raw: 0.0110
\newcommand{\NumBSSolsLo}{0.0087}  % raw: 0.0087
\newcommand{\NumBSSolsHi}{0.0133}  % raw: 0.0133

% docs/RESULTS.md:29, row (2026-10-31 | Report n<U+00B0>3 (descriptive))
%   pattern: `enc-softmax` log-loss ([-+<U+2212>]?(?:\d{1,3}(?:,\d{3})+|\d+)(?:\.\d+)?) \[([-+<U+2212>]?(?:\d{1,3}(?:,\d{3})+|\d+)(?:\.\d+)?), ([-+<U+2212>]?(?:\d{1,3}(?:,\d{3})+|\d+)(?:\.\d+)?)\], BSS ([-+<U+2212>]?(?:\d{1,3}(?:,\d{3})+|\d+)(?:\.\d+)?) \[([-+<U+2212>]?(?:\d{1,3}(?:,\d{3})+|\d+)(?:\.\d+)?), ([-+<U+2212>]?(?:\d{1,3}(?:,\d{3})+|\d+)(?:\.\d+)?)\]
%   matched: `enc-softmax` log-loss 1.0685 [1.0654, 1.0717], BSS 0.0079 [0.0059, 0.0098]
\newcommand{\NumLLencsm}{1.0685}  % raw: 1.0685
\newcommand{\NumLLencsmLo}{1.0654}  % raw: 1.0654
\newcommand{\NumLLencsmHi}{1.0717}  % raw: 1.0717
\newcommand{\NumBSSencsm}{0.0079}  % raw: 0.0079
\newcommand{\NumBSSencsmLo}{0.0059}  % raw: 0.0059
\newcommand{\NumBSSencsmHi}{0.0098}  % raw: 0.0098

% docs/RESULTS.md:29, row (2026-10-31 | Report n<U+00B0>3 (descriptive))
%   pattern: `enc-linear` log-loss ([-+<U+2212>]?(?:\d{1,3}(?:,\d{3})+|\d+)(?:\.\d+)?) \[([-+<U+2212>]?(?:\d{1,3}(?:,\d{3})+|\d+)(?:\.\d+)?), ([-+<U+2212>]?(?:\d{1,3}(?:,\d{3})+|\d+)(?:\.\d+)?)\], BSS ([-+<U+2212>]?(?:\d{1,3}(?:,\d{3})+|\d+)(?:\.\d+)?) \[([-+<U+2212>]?(?:\d{1,3}(?:,\d{3})+|\d+)(?:\.\d+)?), ([-+<U+2212>]?(?:\d{1,3}(?:,\d{3})+|\d+)(?:\.\d+)?)\]
%   matched: `enc-linear` log-loss 1.0728 [1.0699, 1.0756], BSS 0.0030 [0.0016, 0.0044]
\newcommand{\NumLLlin}{1.0728}  % raw: 1.0728
\newcommand{\NumLLlinLo}{1.0699}  % raw: 1.0699
\newcommand{\NumLLlinHi}{1.0756}  % raw: 1.0756
\newcommand{\NumBSSlin}{0.0030}  % raw: 0.0030
\newcommand{\NumBSSlinLo}{0.0016}  % raw: 0.0016
\newcommand{\NumBSSlinHi}{0.0044}  % raw: 0.0044

% docs/RESULTS.md:29, row (2026-10-31 | Report n<U+00B0>3 (descriptive))
%   pattern: `enc-ols`: (\d{1,3}(?:,\d{3})+|\d+) parameters
%   matched: `enc-ols`: 21,003 parameters
\newcommand{\NumParamsOls}{21,003}  % raw: 21,003

% docs/RESULTS.md:29, row (2026-10-31 | Report n<U+00B0>3 (descriptive))
%   pattern: `enc-softmax`: (\d{1,3}(?:,\d{3})+|\d+) parameters
%   matched: `enc-softmax`: 20,995 parameters
\newcommand{\NumParamsSm}{20,995}  % raw: 20,995

% docs/RESULTS.md:29, row (2026-10-31 | Report n<U+00B0>3 (descriptive))
%   pattern: `enc-linear`: (\d{1,3}(?:,\d{3})+|\d+) parameters
%   matched: `enc-linear`: 20,995 parameters
\newcommand{\NumParamsLin}{20,995}  % raw: 20,995

% docs/RESULTS.md:29, row (2026-10-31 | Report n<U+00B0>3 (descriptive))
%   pattern: 6 yearly fits per seed \((\d+) fits\)
%   matched: 6 yearly fits per seed (54 fits)
\newcommand{\NumRepThreeFits}{54}  % raw: 54

% docs/RESULTS.md:29, row (2026-10-31 | Report n<U+00B0>3 (descriptive))
%   pattern: every part verifies and (\d+) files are identical
%   matched: every part verifies and 4 files are identical
\newcommand{\NumRepThreeCheckFiles}{4}  % raw: 4

% docs/HYPOTHESES.md:25, line holding '| H15 |'
%   pattern: population of (\d{1,3}(?:,\d{3})+|\d+) leaks of seed (\d+) \(placement `quantile`\), at H1's (\d+) cuts
%   matched: population of 10,000 leaks of seed 20261118 (placement `quantile`), at H1's 16 cuts
\newcommand{\NumPopN}{10,000}  % raw: 10,000
\newcommand{\NumConfSeed}{20261118}  % raw: 20261118
\newcommand{\NumHoneCutsN}{16}  % raw: 16

% docs/HYPOTHESES.md:25, line holding '| H15 |'
%   pattern: Newcombe 95 % interval \(method (\d+)\)
%   matched: Newcombe 95 % interval (method 10)
\newcommand{\NumNewcombeMethod}{10}  % raw: 10

% docs/HYPOTHESES.md:25, line holding '| H15 |'
%   pattern: against a margin of (\d+) point
%   matched: against a margin of 1 point
\newcommand{\NumHfifteenMargin}{1}  % raw: 1

% docs/HYPOTHESES.md:12, line holding '| H2 |'
%   pattern: with a block-bootstrap (\d+) % interval excluding 0
%   matched: with a block-bootstrap 95 % interval excluding 0
\newcommand{\NumConf}{95}  % raw: 95

% docs/PROTOCOL.md:499, line holding '1. **Confirmatory population.**'
%   pattern: use seed \*\*(\d+)\*\* only
%   matched: use seed **51118** only
\newcommand{\NumDevSeed}{51118}  % raw: 51118

% docs/PROTOCOL.md:500, line holding '2. **Descriptive population.**'
%   pattern: The [\d,]+ leaks of seed (\d+), placement `h1-bands`.*?\(([-+<U+2212>]?(?:\d{1,3}(?:,\d{3})+|\d+)(?:\.\d+)?) % any-cut, ([-+<U+2212>]?(?:\d{1,3}(?:,\d{3})+|\d+)(?:\.\d+)?) % per cut\)
%   matched: The 10,000 leaks of seed 20261029, placement `h1-bands`, as run in iteration 3. Its FI rates are known (97.95 % any-cut, 86.85 % per cut)
\newcommand{\NumDescSeed}{20261029}  % raw: 20261029
\newcommand{\NumDescFIany}{97.95}  % raw: 97.95
\newcommand{\NumDescFIper}{86.85}  % raw: 86.85

% docs/PROTOCOL.md:502, line holding '**The draw, at code level.**'
%   pattern: `level = rng.uniform\(([-+<U+2212>]?(?:\d{1,3}(?:,\d{3})+|\d+)(?:\.\d+)?), ([-+<U+2212>]?(?:\d{1,3}(?:,\d{3})+|\d+)(?:\.\d+)?)\)`
%   matched: `level = rng.uniform(0.02, 0.98)`
\newcommand{\NumLevelLo}{0.02}  % raw: 0.02
\newcommand{\NumLevelHi}{0.98}  % raw: 0.98

% docs/PROTOCOL.md:512, line holding '**Rows:** 502 per pair'
%   pattern: \*\*Rows:\*\* (\d+) per pair
%   matched: **Rows:** 502 per pair
\newcommand{\NumTableRows}{502}  % raw: 502

% docs/PROTOCOL.md:513, line holding '**SHA-256 of `table.json`:**'
%   pattern: `([0-9a-f]{64})`
%   matched: `527e01ba7028fb5666e694fcadd5b22d2fdea31718d96e53d98a29e383183c43`
\newcommand{\NumTableSha}{\texttt{527e01ba7028\ldots{}}}  % raw: 527e01ba7028fb5666e694fcadd5b22d2fdea31718d96e53d98a29e383183c43

% docs/PROTOCOL.md:515, line holding '**Share of levels in [0.02, 0.98]'
%   pattern: body ([-+<U+2212>]?(?:\d{1,3}(?:,\d{3})+|\d+)(?:\.\d+)?) %, upper wick ([-+<U+2212>]?(?:\d{1,3}(?:,\d{3})+|\d+)(?:\.\d+)?) %, lower wick ([-+<U+2212>]?(?:\d{1,3}(?:,\d{3})+|\d+)(?:\.\d+)?) %, range ([-+<U+2212>]?(?:\d{1,3}(?:,\d{3})+|\d+)(?:\.\d+)?) %, fill ([-+<U+2212>]?(?:\d{1,3}(?:,\d{3})+|\d+)(?:\.\d+)?) %, close location ([-+<U+2212>]?(?:\d{1,3}(?:,\d{3})+|\d+)(?:\.\d+)?) %\. Under `h1-bands` each is ([-+<U+2212>]?(?:\d{1,3}(?:,\d{3})+|\d+)(?:\.\d+)?) %
%   matched: body 84.2 %, upper wick 39.7 %, lower wick 31.5 %, range 100 %, fill 67.5 %, close location 55.4 %. Under `h1-bands` each is 50 %
\newcommand{\NumBandBody}{84.2}  % raw: 84.2
\newcommand{\NumBandUpper}{39.7}  % raw: 39.7
\newcommand{\NumBandLower}{31.5}  % raw: 31.5
\newcommand{\NumBandRange}{100}  % raw: 100
\newcommand{\NumBandFill}{67.5}  % raw: 67.5
\newcommand{\NumBandCloseLoc}{55.4}  % raw: 55.4
\newcommand{\NumBandHalf}{50}  % raw: 50

% docs/PROTOCOL.md:516, line holding '6. **What this means for the reading of FI's rate.**'
%   pattern: \((\d{1,3}(?:,\d{3})+|\d+) of (\d{1,3}(?:,\d{3})+|\d+) on a neutral seed\)
%   matched: (8,258 of 10,000 on a neutral seed)
\newcommand{\NumInvSubset}{8,258}  % raw: 8,258
\newcommand{\NumInvOf}{10,000}  % raw: 10,000

% docs/PROTOCOL.md:516, line holding '6. **What this means for the reading of FI's rate.**'
%   pattern: (\d{1,3}(?:,\d{3})+|\d+) of [\d,]+ leaks are tagged `inside`, only (\d+) of them thresholds
%   matched: 4,815 of 10,000 leaks are tagged `inside`, only 520 of them thresholds
\newcommand{\NumInsideTagged}{4,815}  % raw: 4,815
\newcommand{\NumInsideThresholds}{520}  % raw: 520

% docs/PROTOCOL.md:526, line holding '**Shared ground.**'
%   pattern: (\d+) dates, the same row indices on both pairs \((\d+), (\d+), (\d+), (\d+), (\d+), (\d+), (\d+), (\d+)\), chosen by `invariance.choose_cuts` with seed (\d+)
%   matched: 8 dates, the same row indices on both pairs (502, 1048, 1233, 1597, 1994, 2100, 2143, 2664), chosen by `invariance.choose_cuts` with seed 20261027
\newcommand{\NumCutDates}{8}  % raw: 8
\newcommand{\NumCutRowA}{502}  % raw: 502
\newcommand{\NumCutRowB}{1048}  % raw: 1048
\newcommand{\NumCutRowC}{1233}  % raw: 1233
\newcommand{\NumCutRowD}{1597}  % raw: 1597
\newcommand{\NumCutRowE}{1994}  % raw: 1994
\newcommand{\NumCutRowF}{2100}  % raw: 2100
\newcommand{\NumCutRowG}{2143}  % raw: 2143
\newcommand{\NumCutRowH}{2664}  % raw: 2664
\newcommand{\NumCutSeed}{20261027}  % raw: 20261027

% docs/PROTOCOL.md:492, line holding '2. **Compared rows.**'
%   pattern: \(`FUTURE_REACH = (\d+)`\)
%   matched: (`FUTURE_REACH = 10`)
\newcommand{\NumReach}{10}  % raw: 10

% docs/PROTOCOL.md:491, line holding '1. **A missing bar.**'
%   pattern: counts a move above (1e-9)
%   matched: counts a move above 1e-9
\newcommand{\NumCatchThreshold}{\ensuremath{10^{-9}}}  % raw: 1e-9

% docs/PROTOCOL.md:528, line holding '1. **FI (future invariance).**'
%   pattern: <U+00D7>(\d+) and <U+00F7>(\d+) at each cut.*?\(perturbation seed (\d+)\)
%   matched: <U+00D7>3 and <U+00F7>3 at each cut, exactly as `bench/invariance.py` and the adversary run it (perturbation seed 7)
\newcommand{\NumFIup}{3}  % raw: 3
\newcommand{\NumFIdown}{3}  % raw: 3
\newcommand{\NumFIseed}{7}  % raw: 7

% docs/PROTOCOL.md:533, line holding '*Deviation from v1, dated 2026-11-18:*'
%   pattern: future reach 1 \(([-+<U+2212>]?(?:\d{1,3}(?:,\d{3})+|\d+)(?:\.\d+)?) % of draws\)
%   matched: future reach 1 (26.2 % of draws)
\newcommand{\NumReachOneShare}{26.2}  % raw: 26.2

% docs/PROTOCOL.md:535, line holding 'Tile lengths <U+03C4> <U+2208>'
%   pattern: <U+03C4> <U+2208> \{(\d+), (\d+)\} bars
%   matched: <U+03C4> <U+2208> {256, 365} bars
\newcommand{\NumTauA}{256}  % raw: 256
\newcommand{\NumTauB}{365}  % raw: 365

% docs/PROTOCOL.md:535, line holding 'Tile lengths <U+03C4> <U+2208>'
%   pattern: \*\*(\d+) ends in all\.\*\*
%   matched: **28 ends in all.**
\newcommand{\NumTileEnds}{28}  % raw: 28

% docs/PROTOCOL.md:536, line holding 'The deepest past read is'
%   pattern: The deepest past read is (\d+) bars
%   matched: The deepest past read is 189 bars
\newcommand{\NumDeepestRead}{189}  % raw: 189

% docs/PROTOCOL.md:538, line holding '5. **TWO-RUN (ours, after Fonseca's <U+00A7>8.3 class).**'
%   pattern: the bars `t \+ 1 \.\. t \+ (\d+)` \(the only later bars
%   matched: the bars `t + 1 .. t + 10` (the only later bars
\newcommand{\NumTwoRunBars}{10}  % raw: 10

% docs/PROTOCOL.md:539, line holding '**Draws:** `rng = np.random.default_rng('
%   pattern: Development seed \*\*(\d+)\*\* in place of (\d+)\.
%   matched: Development seed **51119** in place of 20261119.
\newcommand{\NumTwoRunDevSeed}{51119}  % raw: 51119
\newcommand{\NumTwoRunSeed}{20261119}  % raw: 20261119

% docs/PROTOCOL.md:540, line holding '**The two <U+03C3> share `z`**'
%   pattern: Primary <U+03C3> = ([-+<U+2212>]?(?:\d{1,3}(?:,\d{3})+|\d+)(?:\.\d+)?); secondary <U+03C3> = ([-+<U+2212>]?(?:\d{1,3}(?:,\d{3})+|\d+)(?:\.\d+)?)\.
%   matched: Primary <U+03C3> = 0.1; secondary <U+03C3> = 1.0.
\newcommand{\NumSigmaLo}{0.1}  % raw: 0.1
\newcommand{\NumSigmaHi}{1.0}  % raw: 1.0

% docs/PROTOCOL.md:547, line holding '**Reported per <U+03C3>, on the perturbed bars:**'
%   pattern: median body ([-+<U+2212>]?(?:\d{1,3}(?:,\d{3})+|\d+)(?:\.\d+)?) to ([-+<U+2212>]?(?:\d{1,3}(?:,\d{3})+|\d+)(?:\.\d+)?), median fill ([-+<U+2212>]?(?:\d{1,3}(?:,\d{3})+|\d+)(?:\.\d+)?) to ([-+<U+2212>]?(?:\d{1,3}(?:,\d{3})+|\d+)(?:\.\d+)?); <U+03C3> = [\d.]+: median body ([-+<U+2212>]?(?:\d{1,3}(?:,\d{3})+|\d+)(?:\.\d+)?)
%   matched: median body 0.019 to 0.096, median fill 0.46 to 0.81; <U+03C3> = 1.0: median body 0.93
\newcommand{\NumBodyBefore}{0.019}  % raw: 0.019
\newcommand{\NumBodyAfter}{0.096}  % raw: 0.096
\newcommand{\NumFillBefore}{0.46}  % raw: 0.46
\newcommand{\NumFillAfter}{0.81}  % raw: 0.81
\newcommand{\NumBodyHiSigma}{0.93}  % raw: 0.93

% docs/PROTOCOL.md:548, line holding '**Runs:** 2 runs per cut in the primary analysis'
%   pattern: reported also with (\d+) run \(run 1\) and (\d+) runs
%   matched: reported also with 1 run (run 1) and 8 runs
\newcommand{\NumTwoRunOne}{1}  % raw: 1
\newcommand{\NumTwoRunMany}{8}  % raw: 8

% docs/PROTOCOL.md:549, line holding '6. **SHUFFLE (the future-shuffle twins'
%   pattern: Reported also with its first (\d+) and (\d+) twins
%   matched: Reported also with its first 1 and 2 twins
\newcommand{\NumShuffleOne}{1}  % raw: 1
\newcommand{\NumShuffleTwo}{2}  % raw: 2

% docs/PROTOCOL.md:495, line holding '5. **The generator's ties to H1.**'
%   pattern: the (\d+)-bar reference window
%   matched: the 90-bar reference window
\newcommand{\NumRefWindow}{90}  % raw: 90

% docs/PROTOCOL.md:510, line holding '**Consequence, added to the generator's "never produced" list'
%   pattern: beyond the (\d+)nd or the (\d+)th percentile
%   matched: beyond the 2nd or the 98th percentile
\newcommand{\NumNeverLo}{2}  % raw: 2
\newcommand{\NumNeverHi}{98}  % raw: 98

% docs/PROTOCOL.md:602, line holding '15. **The cut-date bootstrap**'
%   pattern: the ([\d.]+)th and ([\d.]+)th percentiles of the any-cut difference
%   matched: the 2.5th and 97.5th percentiles of the any-cut difference
\newcommand{\NumBootPLo}{2.5}  % raw: 2.5
\newcommand{\NumBootPHi}{97.5}  % raw: 97.5

% docs/PROTOCOL.md:1097, line holding '2. **The descriptive run:** `python -m olsprime.bench.detectors run'
%   pattern: \*\*Seeds (\d+) \(TWO-RUN\) and (\d+) \(cut-date bootstrap\) are new\*\*
%   matched: **Seeds 29119 (TWO-RUN) and 29120 (cut-date bootstrap) are new**
\newcommand{\NumDescTwoRunSeed}{29119}  % raw: 29119
\newcommand{\NumDescBootSeed}{29120}  % raw: 29120

% docs/PROTOCOL.md:1110, line holding 'takes its correlation estimate <U+03C8> as Fagerland, Lydersen and Laake'
%   pattern: at most ([-+<U+2212>]?(?:\d{1,3}(?:,\d{3})+|\d+)(?:\.\d+)?) point at n = [\d,]+, ([-+<U+2212>]?(?:\d{1,3}(?:,\d{3})+|\d+)(?:\.\d+)?) on an extreme table
%   matched: at most 0.001 point at n = 10,000, 0.003 on an extreme table
\newcommand{\NumPsiDiff}{0.001}  % raw: 0.001
\newcommand{\NumPsiDiffExtreme}{0.003}  % raw: 0.003

% docs/PROTOCOL.md:1110, line holding 'takes its correlation estimate <U+03C8> as Fagerland, Lydersen and Laake'
%   pattern: If a limit falls within ([-+<U+2212>]?(?:\d{1,3}(?:,\d{3})+|\d+)(?:\.\d+)?) point of 0 or of
%   matched: If a limit falls within 0.003 point of 0 or of
\newcommand{\NumPsiFlag}{0.003}  % raw: 0.003

% docs/PROTOCOL.md:1089, line holding '**FI on M2 is three-valued.**'
%   pattern: `high\[t\+1\] <U+2265> (\d+)/(\d+) <U+00B7> close\[t\]`
%   matched: `high[t+1] <U+2265> 2/3 <U+00B7> close[t]`
\newcommand{\NumMtwoNum}{2}  % raw: 2
\newcommand{\NumMtwoDen}{3}  % raw: 3

% docs/PROTOCOL.md:146, line holding '**Decision rules.** A comparison passes when'
%   pattern: gives p < ([\d.]+)
%   matched: gives p < 0.05
\newcommand{\NumPassP}{0.05}  % raw: 0.05

% docs/RESULTS.md:19, row (2026-10-29 | H12)
%   pattern: bootstrap (\d+)th percentile
%   matched: bootstrap 95th percentile
\newcommand{\NumQuantPct}{95}  % raw: 95

% docs/PROTOCOL.md:1123, line holding 'The frozen commit is `b71ba654ff7060ff7f835f8c3325cbd073eddaa9`'
%   pattern: CPython ([\d.]+), numpy ([\d.]+), the development digest
%   matched: CPython 3.12.10, numpy 2.5.4, the development digest
\newcommand{\NumWaveCPython}{3.12.10}  % raw: 3.12.10
\newcommand{\NumWaveNumpy}{2.5.4}  % raw: 2.5.4

% docs/PROTOCOL.md:549, line holding '6. **SHUFFLE (the future-shuffle twins'
%   pattern: K = (\d+) twins per cut, seed (\d+)
%   matched: K = 8 twins per cut, seed 20261030
\newcommand{\NumShuffleK}{8}  % raw: 8
\newcommand{\NumShuffleSeed}{20261030}  % raw: 20261030

% docs/PROTOCOL.md:550, line holding '7. **Run counts per cut, primary:**'
%   pattern: FI (\d+), TRUNC-0 (\d+), TRUNC-1 (\d+), TWO-RUN (\d+), SHUFFLE (\d+)\.
%   matched: FI 2, TRUNC-0 1, TRUNC-1 1, TWO-RUN 2, SHUFFLE 8.
\newcommand{\NumRunsFI}{2}  % raw: 2
\newcommand{\NumRunsTruncZero}{1}  % raw: 1
\newcommand{\NumRunsTruncOne}{1}  % raw: 1
\newcommand{\NumRunsTwo}{2}  % raw: 2
\newcommand{\NumRunsShuffle}{8}  % raw: 8

% docs/PROTOCOL.md:559, line holding '**Secondary schedule (Fonseca's sampled epochs).**'
%   pattern: FI \(perturbation seed (\d+)\).*?(\d+) cuts per pair, \d+ dates shared by both pairs, long cut seed (\d+).*?over the 16 \+ (\d+) cuts
%   matched: FI (perturbation seed 11), TRUNC-0, TRUNC-1 and TWO-RUN (`schedule_id` 2) also run on the adversary's long schedule: 64 cuts per pair, 64 dates shared by both pairs, long cut seed 20261030. SHUFFLE has no seeds there and does not run. "The union over both schedules" is the union of those four detectors over the 16 + 128 cuts
\newcommand{\NumLongFIseed}{11}  % raw: 11
\newcommand{\NumLongPerPair}{64}  % raw: 64
\newcommand{\NumLongCutSeed}{20261030}  % raw: 20261030
\newcommand{\NumLongTotal}{128}  % raw: 128

% docs/PROTOCOL.md:569, line holding 'Leaks with no `At` or `Agg` future node'
%   pattern: about ([-+<U+2212>]?(?:\d{1,3}(?:,\d{3})+|\d+)(?:\.\d+)?) % in the synthetic draw
%   matched: about 9.3 % in the synthetic draw
\newcommand{\NumNoControlShare}{9.3}  % raw: 9.3

% docs/PROTOCOL.md:602, line holding '15. **The cut-date bootstrap**'
%   pattern: (\d{1,3}(?:,\d{3})+|\d+) subsets of (\d+) dates drawn without replacement, the same dates on both pairs, seed (\d+)
%   matched: 2,000 subsets of 8 dates drawn without replacement, the same dates on both pairs, seed 20261120
\newcommand{\NumBootSubsets}{2,000}  % raw: 2,000
\newcommand{\NumBootDates}{8}  % raw: 8
\newcommand{\NumBootSeed}{20261120}  % raw: 20261120

% docs/PROTOCOL.md:617, line holding '**Power.** With 3 % of leaks discordant'
%   pattern: With ([-+<U+2212>]?(?:\d{1,3}(?:,\d{3})+|\d+)(?:\.\d+)?) % of leaks discordant, the paired difference's 95 % half-width is near ([-+<U+2212>]?(?:\d{1,3}(?:,\d{3})+|\d+)(?:\.\d+)?) point and a difference of ([-+<U+2212>]?(?:\d{1,3}(?:,\d{3})+|\d+)(?:\.\d+)?) point is detected with 80 % power; with ([-+<U+2212>]?(?:\d{1,3}(?:,\d{3})+|\d+)(?:\.\d+)?) % discordant, ([-+<U+2212>]?(?:\d{1,3}(?:,\d{3})+|\d+)(?:\.\d+)?) and ([-+<U+2212>]?(?:\d{1,3}(?:,\d{3})+|\d+)(?:\.\d+)?) point
%   matched: With 3 % of leaks discordant, the paired difference's 95 % half-width is near 0.34 point and a difference of 0.49 point is detected with 80 % power; with 8 % discordant, 0.55 and 0.79 point
\newcommand{\NumPowDiscA}{3}  % raw: 3
\newcommand{\NumPowHalfA}{0.34}  % raw: 0.34
\newcommand{\NumPowDetA}{0.49}  % raw: 0.49
\newcommand{\NumPowDiscB}{8}  % raw: 8
\newcommand{\NumPowHalfB}{0.55}  % raw: 0.55
\newcommand{\NumPowDetB}{0.79}  % raw: 0.79

% docs/PROTOCOL.md:616, line holding '**Reported beside it, never deciding:**'
%   pattern: changed the any-cut verdict of (\d+) of (\d{1,3}(?:,\d{3})+|\d+) leaks
%   matched: changed the any-cut verdict of 57 of 4,000 leaks
\newcommand{\NumShiftChanged}{57}  % raw: 57
\newcommand{\NumShiftOf}{4,000}  % raw: 4,000

% docs/PROTOCOL.md:606, line holding '**Equivalence gates** (pass 5a, m4)'
%   pattern: on (\d+) leaks \(equivalence seed (\d+)\)
%   matched: on 50 leaks (equivalence seed 3)
\newcommand{\NumGateLeaks}{50}  % raw: 50
\newcommand{\NumGateSeed}{3}  % raw: 3

% docs/PROTOCOL.md:1076, line holding '- 2026-11-18, before any detector of iteration 5 runs on a leak'
%   pattern: \(register with fixes: (\d+) Blockers, (\d+) Major, (\d+) Minor, (\d+) Nits\) and pass 5a-bis \(register with fixes: no Blocker, (\d+) Major, (\d+) Minor\)
%   matched: (register with fixes: 2 Blockers, 10 Major, 11 Minor, 4 Nits) and pass 5a-bis (register with fixes: no Blocker, 7 Major, 11 Minor)
\newcommand{\NumAfiveaBlock}{2}  % raw: 2
\newcommand{\NumAfiveaMajor}{10}  % raw: 10
\newcommand{\NumAfiveaMinor}{11}  % raw: 11
\newcommand{\NumAfiveaNit}{4}  % raw: 4
\newcommand{\NumAfiveabisMajor}{7}  % raw: 7
\newcommand{\NumAfiveabisMinor}{11}  % raw: 11

% docs/PROTOCOL.md:1086, line holding '- 2026-11-19, after audit pass 5b (GO with fixes'
%   pattern: \(GO with fixes, no Blocker; (\d+) Major, (\d+) Minor, (\d+) Nits\)
%   matched: (GO with fixes, no Blocker; 4 Major, 12 Minor, 5 Nits)
\newcommand{\NumAfivebMajor}{4}  % raw: 4
\newcommand{\NumAfivebMinor}{12}  % raw: 12
\newcommand{\NumAfivebNit}{5}  % raw: 5

% docs/PROTOCOL.md:1122, line holding '- 2026-11-19, after audit pass 5b-bis (GO with fixes'
%   pattern: \(GO with fixes, no Blocker; (\d+) Major, (\d+) Minor, (\d+) Nits;
%   matched: (GO with fixes, no Blocker; 1 Major, 3 Minor, 9 Nits;
\newcommand{\NumAfivebbisMajor}{1}  % raw: 1
\newcommand{\NumAfivebbisMinor}{3}  % raw: 3
\newcommand{\NumAfivebbisNit}{9}  % raw: 9

% docs/PROTOCOL.md:1085, line holding '**The development population's digest (<U+00A7>1.7).**'
%   pattern: SHA-256 `([0-9a-f]{64})`, numpy ([\d.]+),
%   matched: SHA-256 `e03f78b187c1f7a7b62681a6c16214f7bd79d7f13b58459521cacb731ad9bc48`, numpy 2.5.4,
\newcommand{\NumDevDigest}{\texttt{e03f78b187c1\ldots{}}}  % raw: e03f78b187c1f7a7b62681a6c16214f7bd79d7f13b58459521cacb731ad9bc48
\newcommand{\NumNumpyVer}{2.5.4}  % raw: 2.5.4

% docs/PROTOCOL.md:1124, line holding '**Pass 5b-bis reproduced it at the frozen commit:**'
%   pattern: (\d{1,3}(?:,\d{3})+|\d+) passed and (\d+) skipped; the mutant test, (\d+) of (\d+) killed
%   matched: 1,743 passed and 5 skipped; the mutant test, 16 of 16 killed
\newcommand{\NumSuitePassed}{1,743}  % raw: 1,743
\newcommand{\NumSuiteSkipped}{5}  % raw: 5
\newcommand{\NumMutKilled}{16}  % raw: 16
\newcommand{\NumMutTotal}{16}  % raw: 16

% docs/PROTOCOL.md:1125, line holding '**The code digest** (pass 5b-bis, Minor 2).'
%   pattern: `([0-9a-f]{64})` \((\d+) files\)
%   matched: `3b59007a8111cc858bf01e15d7ed6636775903cbabce484ea3ac3b483ab4d5e5` (88 files)
\newcommand{\NumCodeDigest}{\texttt{3b59007a8111\ldots{}}}  % raw: 3b59007a8111cc858bf01e15d7ed6636775903cbabce484ea3ac3b483ab4d5e5
\newcommand{\NumCodeFiles}{88}  % raw: 88

% docs/PROTOCOL.md:1129, line holding '**Minor 3.** Five mutants planted'
%   pattern: the kill count of all (\d+) is reported
%   matched: the kill count of all 21 is reported
\newcommand{\NumMutLater}{21}  % raw: 21

% docs/PROTOCOL.md:1123, line holding 'The frozen commit is `b71ba654ff7060ff7f835f8c3325cbd073eddaa9`'
%   pattern: The frozen commit is `([0-9a-f]{40})`
%   matched: The frozen commit is `b71ba654ff7060ff7f835f8c3325cbd073eddaa9`
\newcommand{\NumFrozenCommit}{\texttt{b71ba654ff70\ldots{}}}  % raw: b71ba654ff7060ff7f835f8c3325cbd073eddaa9

% docs/PROTOCOL.md:720, line holding '**What was read.** To prove the holes D1 and D2'
%   pattern: the (\d+) bars of the 2026-03 raw archive
%   matched: the 31 bars of the 2026-03 raw archive
\newcommand{\NumExposedBars}{31}  % raw: 31

% docs/PROTOCOL.md:97, line holding 'The windows are 30 days, 365 days and all history.'
%   pattern: The windows are (\d+) days, (\d+) days and all history
%   matched: The windows are 30 days, 365 days and all history
\newcommand{\NumWinShort}{30}  % raw: 30
\newcommand{\NumWinYear}{365}  % raw: 365

% docs/PROTOCOL.md:88, line holding '- Features: the 15 `FEATURE_COLUMNS`'
%   pattern: the (\d+) `FEATURE_COLUMNS`
%   matched: the 15 `FEATURE_COLUMNS`
\newcommand{\NumFeatures}{15}  % raw: 15

% docs/PROTOCOL.md:89, line holding '- Labels: three classes over h = 24 bars'
%   pattern: three classes over h = (\d+) bars
%   matched: three classes over h = 24 bars
\newcommand{\NumLabelBars}{24}  % raw: 24

% docs/PROTOCOL.md:83, line holding '233 of the 274 members have a member month in the development period'
%   pattern: of the (\d+) members have a member month.*?the other (\d+) first join
%   matched: of the 274 members have a member month in the development period; the other 41 first join
\newcommand{\NumPanelAll}{274}  % raw: 274
\newcommand{\NumPanelLate}{41}  % raw: 41

% docs/PROTOCOL.md:189, line holding '**Inputs: a window of bars.**'
%   pattern: reads the (\d+) hourly bars of pair s
%   matched: reads the 64 hourly bars of pair s
\newcommand{\NumEncT}{64}  % raw: 64

% docs/PROTOCOL.md:200, line holding '- a linear map of the eight quantities to d_model'
%   pattern: d_model = (\d+)
%   matched: d_model = 32
\newcommand{\NumEncDmodel}{32}  % raw: 32

% docs/PROTOCOL.md:201, line holding 'pre-LayerNorm blocks:'
%   pattern: - (\d+) pre-LayerNorm blocks:
%   matched: - 2 pre-LayerNorm blocks:
\newcommand{\NumEncBlocks}{2}  % raw: 2

% docs/PROTOCOL.md:202, line holding '  - attention, with 4 heads of d_head'
%   pattern: with (\d+) heads of d_head = (\d+);
%   matched: with 4 heads of d_head = 8;
\newcommand{\NumEncHeads}{4}  % raw: 4
\newcommand{\NumEncDhead}{8}  % raw: 8

% docs/PROTOCOL.md:222, line holding '- A readout at decision time d uses the entries with d <U+2212> 30 days'
%   pattern: d <U+2212> (\d+) days < `label_known_at` <U+2264> d
%   matched: d <U+2212> 30 days < `label_known_at` <U+2264> d
\newcommand{\NumEncMemDays}{30}  % raw: 30



% --- draft markers --------------------------------------------------------------------------
\DeclareRobustCommand{\todo}[1]{\textcolor{red}{\textbf{[TODO: #1]}}}
\DeclareRobustCommand{\pend}{\todo{after audit 2a}}
\DeclareRobustCommand{\emptycell}{\textcolor{red}{$\cdot$}}
\newcommand{\bibgap}[1]{\textcolor{red}{[#1]}}
% A method named without a citation: no reading note records it as read (draft rule).
\newcommand{\uncited}[1]{#1\textsuperscript{\textcolor{orange}{[cite?]}}}

\newcommand{\code}[1]{\texttt{#1}}
\newcommand{\FI}{\textsc{fi}}
\newcommand{\TZ}{\textsc{trunc-0}}
\newcommand{\TO}{\textsc{trunc-1}}
\newcommand{\RT}{\textsc{retile}}
\newcommand{\TR}{\textsc{two-run}}
\newcommand{\SH}{\textsc{shuffle}}
\newcommand{\UN}{\textsc{union}}
\newcommand{\pct}{\%}

\title{Measuring a look-ahead detector: future invariance against its rivals on a generated
population of temporal leaks, with a pre-registered case study of OLS attention on crypto bars\\[4pt]
\large\textcolor{red}{[Working title of \code{paper.md}, a draft: Abdou decides the title.]}}
\author{\todo{author name: Abdou's decision}}
\date{Draft 1, 2026-11-19}

\begin{document}
\maketitle

\begin{abstract}
\todo{after audit 2a. The abstract states H15's two labels and the detectors' shares only once
audit phase 2a has recomputed them; until then it is left empty on purpose.}
\end{abstract}

\noindent\textbf{Open decisions (Abdou's, not the draft's):} the title; the author name;
the arXiv category \todo{q-fin.ST with cs.LG, or the reverse}; the licence
\todo{CC BY-NC-SA 4.0 is what the Binance Vision dataset terms impose on derived works; pending,
see \cref{sec:data}}; the endorser \todo{an author with no institution needs an endorsement
from an established arXiv author of the same domain}; any public release of the leak corpus.

% =============================================================================================
\section{Introduction}\label{sec:intro}

A backtest that reads the future overstates its skill, and the leak rarely announces itself: a
label read one bar early, a scaler fitted on the test month, a comparison with the next bar's
close. Several tools test for such leaks by running a pipeline twice, once on the data as they
are and once with the data after a cut changed or removed, and by flagging any output at or
before the cut that moves. Fonseca's two-run differential detector perturbs the future at sampled
epochs and flags a change~\cite{fonseca2026}. % section 8.3 (read; re-read verbatim 2026-11-18, PROTOCOL 0.3)
Freqtrade's lookahead analysis ends a backtest one candle after a trade and compares the indicators
of that truncated run with those of the full run~\cite{freqtrade}. % lookahead.py, read verbatim (PROTOCOL 0.4); docs page
Causify DataFlow re-runs a computation on tiles of the history and compares the result with the
full run up to numerical precision~\cite{causify2025}. % Remark 4.25 and Theorem 4.24 (read; verbatim per PROTOCOL 0.4)

What we did not find, in the works we read, is a \emph{measurement} of such detectors: how many
leaks of a stated population each one catches, with an interval; whether it ever flags a
function that does not read the future; and why it misses what it misses. Fonseca evaluates a
static checker and two dynamic detectors on a fixed corpus of pipelines written in the paper's own
language, and states that such a corpus shows that leak types exist, not how often they
occur~\cite{fonseca2026}. % section 8 (corpus, Table 2) and section 9 (threats to validity), read
This paper reports a measurement of that kind for one bench, with every choice registered before
the population it is measured on was drawn.

Our own detector is \emph{future invariance} (\FI): at a cut $t$, every value of every bar after
$t$ is multiplied by \NumFIup{} in one run and divided by \NumFIdown{} in the other, the next bar
being built so that its comparisons with the recent past flip in exactly one of the two runs, and
every output at or before $t$ must stay identical (\cref{sec:fi}). \FI{} is not a new kind of
test: it is an instance of Fonseca's two-run class~\cite{fonseca2026}, % section 8.3
with its perturbation, cut schedule and catch rule fixed by us. We measure it against five
rivals on one population of \NumPopN{} generated leaks, each paired with an honest twin that reads
the past only: truncation at the cut (\TZ); freqtrade's comparison, rebuilt on our expressions
(\TO); Causify's re-tiling (\RT); a two-run detector with random perturbations, defined by us
because the literature we read gives no default (\TR); and a detector that replaces the future by
bars drawn from the past (\SH). The detectors, the metrics, the denominators and the single
deciding comparison, H15, were registered and audited before any detector touched the population.
H15's prior was that truncation catches at least as many leaks as \FI{}: the measurement had to be
worth reporting if our detector lost it.

\pend{} \emph{[One paragraph, written after 2a: H15's sign and margin labels, the shares of the
detectors under the registered denominators, and the registered sentences of
\cref{sec:prereg-claims}. It must read the same whichever detector the numbers favour.]}

\paragraph{Contributions, in the order of their strength.}
\begin{enumerate}[leftmargin=*]
\item \textbf{The measurement} (\cref{sec:bench,sec:prereg,sec:results}). Catch rates with Wilson
  intervals, honest-twin false positives and a classification of misses, for six detectors and
  their union on the same fresh population, under three conventions for reading a missing bar,
  with a pre-registered comparison and its dated changes.
\item \textbf{The size and power of a pass rule} for paired forecast comparisons on overlapping
  hourly labels, measured by block sign-flips (\cref{sec:power}). It turns one of our negative
  verdicts into what it is: a difference too small for the rule to see.
\item \textbf{Matrix} (\cref{sec:matrix}), the memory our forecasts read: an append-only log of
  (features, outcome, time) entries, never compacted, with a binary tree of sufficient statistics
  over it, from which the ridge forecast of any window of entries is computed at query time and
  split into the contributions of sub-periods and of features. We claim no priority; we state what
  we did not find in the parts of the works we read.
\item \textbf{A case study} (\cref{sec:case}): pre-registered forecasting hypotheses on hourly
  crypto bars, reported exactly as our results log records them, failures and an invalid
  measurement included.
\end{enumerate}

\paragraph{What we do not claim.} No profit, and no trading result of any kind. No result from the
validation year or from the vault, both sealed (\cref{sec:limits} states the one exposure of
vault rows, in an early audit). No rate over ``all leaks'': every rate is this generator's, on
these bars. No ranking of \FI{} against truncation in general. No claim that the two-run idea is
ours. No claim that OLS-based attention forecasts better than softmax attention or is cheaper to
run: the registered test is one of non-inferiority, and the cost measurement was invalid
(\cref{sec:case}).

% =============================================================================================
\section{Related work}\label{sec:related}

\paragraph{How the sources were read.} Every source below was read in part, and through a fetch
tool that summarises web pages unless a page is said to have been read as an image or verbatim.
Each citation is used only for a statement found in the parts our reading notes record as read;
the part is named beside it in the source of this draft. \todo{The orchestrator re-reads each cited
passage verbatim before submission.}

\paragraph{Look-ahead detectors.}
Fonseca~\cite{fonseca2026} % sections 8.2 (end), 8.3, 9; Table 2. Read in body text; 1-4 not read.
compares a static checker, an effect system over the paper's language, with dynamic detectors.
Its two-run differential detector re-executes the pipeline on perturbed future data at sampled
epochs and flags a change; its tiling detector compares outputs across shifted data windows; a
separate dynamic oracle recomputes a decision after perturbing every datum whose availability
exceeds the sampled epoch. % 8.3 (two-run, tiling), end of 8.2 (oracle)
The text we read gives no perturbation law, no rule for the epochs, no number of runs and no
tolerance for the two-run detector, % PROTOCOL 0.3: all 11 occurrences of "perturb" checked verbatim
so there is no default two-run detector to copy: our \TR{} is our own definition. Fonseca names
leak classes by the way a detector misses them: a leak placed between sampled epochs (M1), a future
value read through a threshold that the perturbation does not flip (M2), and a future-valued
series re-stamped to an admissible time (M3); and clean classes: a future value subtracted from
itself (FP1) and an operation that ignores its future argument (FP2), which the checker wrongly
rejects, and an availability that reduces to the present (FP2b), which it accepts. % 8.3 and Table 2; no code example given for any class
We rebuild M1--M3, FP1 and FP2 on bars (\cref{sec:controls}); FP2b needs availability metadata
our bars do not carry. Fonseca's corpus is small and clustered by mechanism, and the paper reads it as
evidence that leak types exist, not of how often they occur; % section 9
we make no comparison with its rates, which come from another corpus in another language.

Freqtrade's lookahead analysis~\cite{freqtrade} % lookahead.py (develop, SHA-256 in refs.bib) and the docs page
runs one full backtest, then, for each analysed signal, a backtest whose data end one candle after
the signal's date, because the strategy does not act on the last candle. It compares every kept
indicator row of the truncated run with the full run's, exactly, with a missing value equal to a
missing value, and reports a bias if any indicator differs or a signal is not reproduced. Its
documentation notes that only signals that fired are verified. It truncates the future; it does
not perturb it.

Causify DataFlow~\cite{causify2025} % Theorem 4.24 and Remark 4.25 (section 4.4); read to Remark 4.30
splits a history into tiles of length at least the computation's context window, re-runs the
computation with two-tile input windows, and checks that each tiling reproduces the full run up to
numerical precision; a mismatch means that the window is too short, that the node reads the
future, or that it is not deterministic. It re-cuts the history; it does not perturb it either.

\paragraph{Least squares as a weighted sum of past outcomes.}
Goulet Coulombe~\cite{gouletcoulombe2025ols} % Eq. (7) p. 5; p. 4; Eq. (23) p. 7; Appendix A.4 pp. 15-16; p. 10
writes the OLS prediction for a test row as a sum, over training rows, of a similarity of
transformed regressors times the training outcome, in the query--key--value notation of attention,
from a fixed training matrix to the test rows; the weights can be negative, and a ridge version is
given. His Appendix~A.4 warns that masking future entries does not remove a leak when the inverse
Gram matrix is estimated on the full sample; our protocol fits every statistic inside the
walk-forward fold for that reason. Goulet Coulombe, G\"obel and Klieber~\cite{gouletcoulombe2024dual}
% section 2.1 and Table 2 (summarised read)
decompose an out-of-sample ridge or kernel forecast into the contributions of the training
observations, through the dual form on a fixed training set, with a separate per-feature route;
the observation weights can be negative. Per-observation decomposition of a ridge forecast is
therefore prior art, and so are negative weights. Garnelo and Czarnecki~\cite{garnelo2023}
% sections 2-3, Def. 3.1, Thm 3.2 (summarised read)
build key--value--query modules whose output is a regularised least-squares solve in closed form,
the regulariser being either a hyperparameter or learned.

\paragraph{Memories of statistics.}
MesaNet~\cite{vonoswald2025mesanet} % Eqs. (4), (6), (8), pp. 1-7 read
defines a layer that solves a ridge regression over the causal prefix of a sequence, keeps running
statistics of the keys and values updated through forget and input gates, and reads the solution
out by conjugate gradient. Log-Linear Attention~\cite{guo2025loglinear} % section 3.1 pp. 4-5, Table 1
partitions the causal prefix with a Fenwick tree into a logarithmic number of segments, each
summarised by a linear-attention state, and mixes their readouts with data-dependent non-negative
weights; buckets are merged as the sequence grows, and the readout is over the whole prefix.
Regression cubes~\cite{chen2006regcube} % pp. 1-10 read as images: Def. 2, Thms 4.1, 4.3, 4.4
compress the regression of a data-cube cell into a measure whose size does not depend on the number
of tuples, and aggregate cells losslessly through sums of their measures, along time as along the
other dimensions, over a lattice of stored cells.

\paragraph{Return prediction.} Nagel~\cite{nagel2025} % sections I to II.F, pp. 1-22
shows that, for ridgeless regressions on many random features over a short rolling window of
monthly returns, the forecast is a similarity-weighted average of past returns that tracks a
volatility-timed momentum rule. We use it only to motivate an exploratory control of our case study
(\cref{sec:case-explore}), not as a result about our bars.

\paragraph{Not cited in this draft.} Statistical methods named in the text (Wilson's interval,
Newcombe's interval for paired proportions, McNemar's test, the Diebold--Mariano test with the
Harvey--Leybourne--Newbold correction, the stationary bootstrap with the Politis--White block
length, Page's CUSUM, the recursive residuals of Brown, Durbin and Evans, the pairwise update of
Chan, Golub and LeVeque, and Fagerland, Lydersen and Laake's implementation of Newcombe's
interval) and the linear attention of Katharopoulos et al.\ carry no citation yet:
no reading note records them. \todo{Read and cite each, or keep them as named methods.} Our design
notes credit the idea of an append-only log with a tree of summaries to OptMem; it was not read
for this draft. \todo{Read OptMem before crediting it in \cref{sec:matrix}.}

% =============================================================================================
\section{The bench}\label{sec:bench}

\subsection{Bars, cuts and the catch rule}\label{sec:bars}
The bars are the daily spot klines of BTCUSDT and ETHUSDT from Binance Vision, read through a seal
that drops every bar of the vault and, for this bench, kept to the development period (August 2017
to December 2024). A \emph{cut} is a row $t$ of a pair's series; the primary schedule is H1's
\NumHoneCutsN{} cuts: \NumCutDates{} dates with the same row indices on both pairs
(\NumCutRowA, \NumCutRowB, \NumCutRowC, \NumCutRowD, \NumCutRowE, \NumCutRowF, \NumCutRowG,
\NumCutRowH), drawn with seed \NumCutSeed. A leak reads at most \NumReach{} bars ahead, so only
rows $t-\NumReach,\dots,t$ can read a bar after $t$; the scored rows among them are the
\emph{compared rows}. A leak is \emph{caught at a cut} if, at some compared row, one of its
outputs moves by more than \NumCatchThreshold, turns missing or stops being missing; it is
\emph{caught} if it is caught at one cut at least. A move at the level of rounding never counts.

\subsection{Future invariance, and its record before this benchmark}\label{sec:fi}
\FI{} multiplies every value of every bar after the cut (prices, volumes, trades, taker volumes)
by \NumFIup{} in one run and divides it by \NumFIdown{} in the other; every ratio inside a bar
moves too, and only the time stamps stay. Bar $t+1$ is built rather than drawn: each of its
quantities sits above the recent bars in one run and below them in the other, and its candle takes
opposite shapes and colours in the two runs, so that a comparison between bar $t+1$ and the recent
past flips in exactly one run, by construction, inside stated bands. Outside those bands, and for
comparisons two or more bars ahead, a flip is likely but not guaranteed. A model that reads the time
stamps passes the test; code review and the seal must stop it.

On a seeded population of \NumAdvN{} random leaks (seed \NumAdvSeed, the population that becomes
the \emph{descriptive} one below), with an honest twin each, \FI{} caught \NumAdvCaught{} of
\NumAdvTotal{} leaks at one of its \NumAdvCuts{} cuts at least, \NumAdvRate\pct{}
(\NumConf\pct{} Wilson interval [\NumAdvRateLo\pct, \NumAdvRateHi\pct]); the mean per-cut rate
was \NumAdvPerCut\pct{} [\NumAdvPerCutLo\pct, \NumAdvPerCutHi\pct]. No honest twin was caught
(\NumAdvTwinsCaught{} of \NumAdvTwinsTotal), and \NumAdvTwinsConst{} twins (\NumAdvTwinsConstPct\pct)
were constant at the compared rows of both schedules, so that nothing could have caught them there. These are this
generator's rates, not a rate over leaks in general. \NumAdvMissed{} leaks were caught at no cut.
How many of them are constant at the compared rows depends on the set of cuts over which constancy
is judged, and our two earlier records used two sets: over all \NumCorrAllCuts{} cuts of the two
schedules, \NumCorrBothConst{} misses are constant, \NumCorrBothUnch{} are unchanged under every
future-shuffle twin and \NumCorrBothBlind{} are true blind spots; over H1's \NumHoneCutsN{} cuts
only, the split is \NumCorrHoneConst{}, \NumCorrHoneUnch{} and \NumCorrHoneBlind. The two sets are
nested, with \NumCorrBetween{} leaks between them, and every count of constant misses in this paper
names its set of cuts. Half of that generator's thresholds were
drawn inside the very bands where \FI{} flips by construction, so these rates are not independent of
\FI's own design; this is why the benchmark's confirmatory population is drawn with another
placement (\cref{sec:generator}).

\subsection{The generator}\label{sec:generator}
A leak is a deterministic function of one pair's bars, with no fit, that gives one output per row
$s$ (a binary target, a continuous value, or three probabilities) and reads at least one bar after
$s$, at most \NumReach{} after it and none more than \NumDeepestRead{} before it. Its quantities
are the bar's levels, the ratios inside a bar (VWAP, trade size, taker shares, body, wicks, range)
and its shapes (fill, close location, wick ratio, colour). The generator draws six families: a
level of a future bar against the past or another future bar; a ratio against the past or a
threshold; candle shapes and colours; aggregates over the next bars; continuous functions of a
future value; and logical or weighted combinations of two leaks. A comparison with a missing value
is false, so a binary leak never outputs a missing value, while a continuous one does. Every leak
has an \emph{honest twin}, the same expression with each future offset moved into the past, which
no detector can move. The generator never produces a fitted leak, a read more than \NumReach{} bars
ahead, or a read of the time stamps, and under the quantile placement below no threshold beyond the
\NumNeverLo{}nd or the \NumNeverHi{}th percentile of its quantity's pre-2019 values.

\paragraph{Two placements of thresholds.} The \emph{descriptive} population (seed \NumDescSeed,
placement \code{h1-bands}) is the one of \cref{sec:fi}: it places half of its thresholds on body,
wicks, range, fill and close location inside \FI's bands. Its \FI{} rates and its misses were seen
before this design, so every number on it is descriptive. The \emph{confirmatory} population (seed
\NumConfSeed, \NumPopN{} leaks, placement \code{quantile}) draws each threshold as a quantile of
the quantity's pre-2019 values on the same two pairs (\NumTableRows{} rows per pair, none of them
scored), at a level uniform on [\NumLevelLo, \NumLevelHi]. The quantile placement follows the bars;
it is not neutral to \FI. Natural thresholds fall inside \FI's band more often than half the time
for some quantities and less often for others (\cref{tab:bands}). Other ties between the generator
and \FI{} remain in both placements: the \code{designed} tag, a \NumRefWindow-bar reference window, the
margin position, the weight of a one-bar future offset, and a body boundary shared with
\FI{} half the time. A \emph{placement-invariant subset}, the leaks that draw no threshold and are
identical under both placements (\NumInvSubset{} of \NumInvOf{} on a neutral seed), is registered
for a sensitivity reading, never a deciding one.

\begin{table}[t]
\centering
\caption{Share of quantile levels in [\NumLevelLo, \NumLevelHi] whose pre-2019 quantile falls
inside \FI's band, per quantity. Under the \code{h1-bands} placement each share is \NumBandHalf\pct{}
by construction. Source: the registered quantile table (SHA-256 \NumTableSha).}
\label{tab:bands}
\begin{tabular}{lcccccc}
\toprule
 & body & upper wick & lower wick & range & fill & close location \\
\midrule
inside the band (\%) & \NumBandBody & \NumBandUpper & \NumBandLower & \NumBandRange & \NumBandFill & \NumBandCloseLoc \\
\bottomrule
\end{tabular}
\end{table}

\subsection{The detectors}\label{sec:detectors}
All detectors of the primary analysis share H1's \NumHoneCutsN{} cuts and its catch rule
(\cref{tab:detectors}). They do not make the same number of runs per cut, so every comparison is
also reported at equal run counts.

\begin{table}[t]
\centering
\caption{The detectors as registered. Runs are re-executions per cut in the primary analysis.}
\label{tab:detectors}
\small
\begin{tabularx}{\linewidth}{>{\raggedright\arraybackslash}p{2.1cm}>{\raggedright\arraybackslash}Xc}
\toprule
Detector & What changes after the cut, and what is compared & Runs \\
\midrule
\FI & \FI's perturbation, $\times\NumFIup$ and $\div\NumFIdown$ (seed \NumFIseed); caught if either
direction moves an output; also reported per direction. & \NumRunsFI \\
\TZ & The input ends at bar $t$: every later bar is missing. & \NumRunsTruncZero \\
\TO & Freqtrade's comparison, rebuilt on our expressions: the input ends at bar $t+1$, and the
compared rows include the slice's last row, $t+1$. It is \TZ{} one bar later: a date-shift
replicate, not a second mechanism. & \NumRunsTruncOne \\
\RT & Causify's protocol: the bars are sliced to two-tile input windows, tiles of
\NumTauA{} and \NumTauB{} bars anchored at the first scored row (\NumTileEnds{} tile ends in all),
and every scored row of a tile is compared with the full run. & --- \\
\TR & Ours, after Fonseca's class: the bars $t+1,\dots,t+\NumTwoRunBars$ get independent
log-normal factors on each column, with repairs that keep each bar valid (wicks scaled in log units,
VWAP clipped to the range, taker volumes capped); scale $\sigma=\NumSigmaLo$ (primary) or
\NumSigmaHi{} (secondary), with common random numbers; also \NumTwoRunOne{} and \NumTwoRunMany{} runs. &
\NumRunsTwo \\
\SH & The future-shuffle twins used as a detector: each bar after the cut keeps its time stamps
and takes the values of a bar drawn with replacement from the same pair's bars at or before the cut;
\NumShuffleK{} twins per cut (seed \NumShuffleSeed), also the first \NumShuffleOne{} and \NumShuffleTwo. & \NumRunsShuffle \\
\UN & Caught by any of \FI, \TZ, \TR{} ($\sigma=\NumSigmaLo$, \NumRunsTwo{} runs) and \SH{}
(\NumShuffleK{} twins). & --- \\
\bottomrule
\end{tabularx}
\end{table}

\TO{} and \RT{} are truncation at other dates. On fit-free functions whose past reach fits in the
input window, \RT's flags equal those of \TZ{} run at the tile's end, and this equality is asserted
by the code; \RT{} is therefore reported as truncation on another schedule, never in the paired
comparisons. As run here, neither truncation nor re-tiling applies to fitted models, whose honest
outputs a shortened input window would change. \TR's law inflates bodies: on synthetic bars, at
$\sigma=\NumSigmaLo$, the median body grew from \NumBodyBefore{} to \NumBodyAfter{} and the median
fill from \NumFillBefore{} to \NumFillAfter; at $\sigma=\NumSigmaHi$ the median body reached
\NumBodyHiSigma. A secondary schedule follows Fonseca's sampled epochs: \FI, \TZ, \TO{} and \TR{}
also run at \NumLongPerPair{} cuts per pair drawn at random (seed \NumLongCutSeed).

\paragraph{Three conventions for a missing read.} Truncation makes a bar missing, and how a
comparison with a missing operand is read decides whether a binary leak's output changes. The
conventions are defined on the atomic comparisons inside a leak and apply to both runs: under G,
the generator's own and the primary, such a comparison is false; under M, the mirror, it is true;
under P, propagation, it is missing and so is every node above it, so that \TZ{} flags a leak at
every cut unless its output was already missing wherever truncation removes a read. \FI, \TR{} and \SH{} keep every bar present
and run under G only.

\subsection{Controls}\label{sec:controls}
The controls have predicted verdicts and are never pooled into any rate.
(i)~Every honest twin, under every detector, convention, schedule and catch rule: any flag is a
bench bug. (ii)~Cancelled controls after Fonseca's FP1 and FP2, generated from each leak: the twin
with its future node $v$ replaced by $w + (v - v)$, by $\min(w, w + 0\cdot v)$, or by the
missing-aware $\mathrm{fmin}(w, w + 0\cdot v)$; their flags are predicted by code that shares the
expression evaluator with the detectors but not their comparison, so we call this a code-level
prediction. About \NumNoControlShare\pct{} of leaks have no future node of the required kind and get
no control. (iii)~Fonseca's classes written by hand on our bars, with their predicted verdicts
committed as tests before any detector ran on real bars (two of these predictions were replaced
after registration, \cref{sec:changes}): M1, a comparison of the next close with the
present one, masked on every compared row; M2, the next close against twice the present one, and
its mirror; M3, the next close frozen and stored as a quantity of the present bar. FP2b is not
rebuilt. (iv)~Assertions on the implementation: every truncation catch goes through a missing read;
under P, \TZ{} flags every leak at every cut unless its output is already missing where the read
disappears; a leak whose root is one comparison and its mirror are flagged together wherever every
read is defined in both runs and the comparison is not tied. These are identities of the code or of
IEEE arithmetic: a failure is a bug, never evidence.

\subsection{Metrics}\label{sec:metrics}
For each detector, on the confirmatory population, under convention G unless stated:
the share caught at one cut at least, with its \uncited{Wilson} interval, also by family and by threshold
position; the mean per-cut rate; the share caught within the first $m$ cuts in date order; the split
of catches into those through a missing read and those through a changed value; the three
conventions; three denominators (all \NumPopN{} leaks, the leaks not constant at the compared rows
of H1's cuts, and the leaks the generator does not mark as covered by \FI's construction); for every
pair among \FI, \TZ, \TR{} and \SH, the paired $2\times 2$ table, the discordant counts, the
difference of shares with \uncited{Newcombe's} interval for paired proportions (method
\NumNewcombeMethod) and \uncited{McNemar's} exact $p$-value; each detector's mirror asymmetry; the
leaks the union catches and each detector misses; and the union's misses, classed by constancy under
both definitions of \cref{sec:fi} and by the secondary arms that catch them. No adjusted score and no
``best detector'' is chosen after the run.

% =============================================================================================
\section{The pre-registration and its dated changes}\label{sec:prereg}

\paragraph{H15.} The one deciding comparison is \TZ{} against \FI, on H1's \NumHoneCutsN{} cuts,
with the any-cut share over all \NumPopN{} leaks of seed \NumConfSeed, under convention G. Its
statistic is Newcombe's \NumConf\pct{} interval $[L, U]$ (method \NumNewcombeMethod) of \TZ's share
minus \FI's, in points, against a smallest difference of interest of \NumHfifteenMargin{} point.
Two labels are reported. The sign: ``above \FI'' if $L>0$, ``below \FI'' if $U<0$, otherwise ``not
separated from \FI''. The margin, when the sign is not ``above'': ``within \NumHfifteenMargin{}
point'' if $L > -\NumHfifteenMargin$, ``by more than \NumHfifteenMargin{} point'' if
$U < -\NumHfifteenMargin$, otherwise ``margin not resolved''. The prior, that truncation catches at
least as many leaks, holds if the sign is ``above''; is not refuted if the sign is ``not separated''
and $L > -\NumHfifteenMargin$; fails if the sign is ``below''; and is undecided otherwise. One
qualifier is registered: if the difference has the opposite sign, with an interval that excludes 0,
under convention M or on the leaks not constant at the compared rows, the verdict is called
``convention-dependent'' and no direction is claimed. A planning computation gave a
\NumConf\pct{} half-width near \NumPowHalfA{} point with \NumPowDiscA\pct{} of leaks discordant and
\NumPowHalfB{} point with \NumPowDiscB\pct{}, so \NumPopN{} leaks suffice for a
\NumHfifteenMargin-point margin. Every interval is conditional on these bars and on
\NumCutDates{} cut dates shared by two pairs.

\paragraph{Beside H15, never deciding.} The mean per-cut rates; convention P; the first-$m$-cuts
curve; the comparisons at equal run counts; the McNemar $p$-value; the placement-invariant subset;
\TO{} against \TZ{} as a one-bar date shift (on synthetic bars it changed the any-cut verdict of
\NumShiftChanged{} of \NumShiftOf{} leaks, of the order of the margin); and a \emph{cut-date
bootstrap}: \TZ{} and \FI{} on the long schedule's \NumLongPerPair{} dates, \NumBootSubsets{} subsets
of \NumBootDates{} dates drawn without replacement (seed \NumBootSeed), whose \NumBootPLo th and
\NumBootPHi th percentiles of the any-cut difference are reported beside H15's interval. If that
spread is wider than the margin, the registered rule makes the paper say that H15's verdict is
date-dependent.

\paragraph{The sequence.} The design was drafted on 2026-11-18 and audited twice before
registration by an adversarial reviewer (\cref{sec:ai}): pass 5a found \NumAfiveaBlock{} blockers,
\NumAfiveaMajor{} major, \NumAfiveaMinor{} minor findings and \NumAfiveaNit{} nits; pass 5a-bis, on
the revised draft, no blocker, \NumAfiveabisMajor{} major and \NumAfiveabisMinor{} minor. Every
finding was fixed or logged. Two of them reshaped the design: the rule that ignores outputs turning
missing does not isolate catches through a missing bar, and truncation's rate depends on how a
comparison with a missing value is read; hence the read split and the three conventions. A draft
hypothesis on why truncation would win was dropped: one half of it held by construction, since every
truncation catch goes through a missing read, and the other was fixed by the generator's family
weights. The registration was dated in the protocol's change log, with the deviations from our own
brief: the primary analysis holds every detector to H1's cuts, the brief's sampled dates becoming
the secondary schedule, and the brief's ``truncation one bar later'' becoming freqtrade's comparison,
whose compared rows were corrected to include row $t+1$ (without it, the arm could not catch a leak
of future reach one, \NumReachOneShare\pct{} of draws).

The only contact with real bars before the official runs, besides reading two earlier outputs for
counts, was the quantile table, computed from the pre-2019 bars on 2026-11-18 and registered by its
digest. The protocol records that it was computed once and that its interpolation method was chosen
before the read; neither can be verified from any commit. The detectors were then built and tested on synthetic bars and on a development seed
(\NumDevSeed) that is never reported. The code was audited (pass 5b: no blocker, \NumAfivebMajor{}
major, \NumAfivebMinor{} minor, \NumAfivebNit{} nits), fixed, and audited again on the diff of the
fixes (pass 5b-bis: no blocker, \NumAfivebbisMajor{} major, \NumAfivebbisMinor{} minor,
\NumAfivebbisNit{} nits). The frozen commit is \code{b71ba65}. At it, the test suite gave
\NumSuitePassed{} passed and \NumSuiteSkipped{} skipped, and an opt-in test killed \NumMutKilled{} of
\NumMutTotal{} planted mutants. Five further mutants, planted by the last audit, survive the
tests; the tests that would kill them are deferred until after the runs, and the kill count over all
\NumMutLater{} will then be reported.

\paragraph{The platform pin.} The population a seed produces depends on the platform's
exponential: \code{exp} and \code{log} differ by one unit in the last place between the Windows C
runtime and glibc, and the same code gives another development digest in WSL2. The official runs
were therefore pinned to one environment (Windows~11, AMD64, CPython \NumWaveCPython, numpy
\NumWaveNumpy), and the run command refuses to read a bar unless it regenerates the development
population with the registered digest (\NumDevDigest).

\paragraph{Changes after registration.}\label{sec:changes} On 2026-11-19, after pass 5b and before
any detector ran on real bars, a dated line of the protocol changed the registered text, and where it
and the earlier text disagree it is normative.
(i)~Two predicted verdicts were replaced. \FI's verdict on M2 became three-valued: at a cut where the
next bar's high is at least \NumMtwoNum/\NumMtwoDen{} of the present close, \FI{} must flag M2; at
every other cut the verdict is ``not predicted'', because the registered half of the rule was wrong
(the run that multiplies the future can lift the next close above twice the present one through the
recent window's high alone). Under convention P, the truncation detectors must flag M1 at every cut,
except a cut whose truncation point is the series' last row, where nothing is removed.
(ii)~Contrary to the first registration, \FI's verdicts on M2 depend on the bars and so are not
identities of the code; a disagreement with them still stops the run.
(iii)~The descriptive run received its own \TR{} and bootstrap seeds, \NumDescTwoRunSeed{} and
\NumDescBootSeed, chosen before any run with them.
(iv)~Narrowings of the implementation, ruled faithful by the audit: Newcombe's correlation estimate
is taken as \uncited{Fagerland, Lydersen and Laake} implement method \NumNewcombeMethod, not from a
re-reading of Newcombe's paper, a difference of at most \NumPsiDiff{} point at \NumPopN{} leaks and
\NumPsiDiffExtreme{} on an extreme table, so the paper must say so if a limit falls within
\NumPsiFlag{} point of 0 or of $-\NumHfifteenMargin$; the cancelled controls do not run through
\RT; the read split is counted at row level; the cut-date spread is the difference of the two
percentiles.
(v)~Metrics were implemented as registered: the sensitivity ``ignores outputs that turn missing'',
and the three denominators applied to every share.

\paragraph{The stop rule and the spent seeds.} Every assertion and gate is evaluated first, and no
rate or per-leak verdict is written until all pass; a failure stops the run, and the next run is
reported as made after a stopped attempt. The official runs were made on 2026-11-19 on Windows, at
the commit that names the frozen one (\code{f12b9b5}), in this order: the registered digest check;
a replay of the earlier adversary on the descriptive seed; the descriptive run; the confirmatory
run, the first and only draw of seed \NumConfSeed{} and of the \TR{} and bootstrap seeds
\NumTwoRunSeed{} and \NumBootSeed; and an identity rerun of it with another number of workers.
These three seeds are spent. Whether every assertion and gate passed, the outputs' digests and the
rerun's byte identity are recorded in the protocol's change log and are being checked by audit
phase 2a: \pend.

\paragraph{What this paper may and may not claim from these runs.}\label{sec:prereg-claims}
These limits were registered with the design. It may state the catch rates of these detectors on this generator and
these bars, with intervals; their overlap; the dependence of truncation's rate on the convention,
in the registered sentence ``truncation's rate lies between $X$ and $Y$ depending on how a missing
bar is read''; each detector's mirror asymmetry; the controls' verdicts; that our two-run arm is our
definition and that its law inflates bodies; and H15's two labels, with the qualifier if triggered
and the cut-date spread beside them. It may not state a rate over all leaks; that \FI{} is better or
worse than truncation in general; that \FI's rate on this population is free of \FI's bands; that
re-tiling is a mechanism independent of truncation here; that truncation or re-tiling as run here
applies to fitted models; a rate of freqtrade or of Causify; or any comparison with Fonseca's
numbers.

% =============================================================================================
\section{Results of the benchmark}\label{sec:results}

\todo{after audit 2a: every number of this section, H15's labels, the qualifier's cells, the
cut-date percentiles and whether the date-dependent sentence fires. The tables below are laid out
with empty cells. If a Newcombe limit falls within \NumPsiFlag{} point of 0 or of
$-\NumHfifteenMargin$, say so (\cref{sec:changes}).}

\begin{table}[t]
\centering
\caption{Any-cut shares on the confirmatory population (seed \NumConfSeed, \NumPopN{} leaks,
H1's \NumHoneCutsN{} cuts; \RT{} at its \NumTileEnds{} tile ends), with \NumConf\pct{} Wilson
intervals, under three denominators. \pend}
\label{tab:shares}
\small
\begin{tabular}{llcccc}
\toprule
Detector & Conv. & Runs/cut & All leaks & Not constant & \code{designed: no} \\
\midrule
\FI & G & \NumRunsFI & \emptycell & \emptycell & \emptycell \\
\TZ & G & \NumRunsTruncZero & \emptycell & \emptycell & \emptycell \\
\TZ & M & \NumRunsTruncZero & \emptycell & \emptycell & \emptycell \\
\TZ & P & \NumRunsTruncZero & \emptycell & \emptycell & \emptycell \\
\TO & G & \NumRunsTruncOne & \emptycell & \emptycell & \emptycell \\
\RT & G & --- & \emptycell & \emptycell & \emptycell \\
\TR{} ($\sigma=\NumSigmaLo$) & G & \NumRunsTwo & \emptycell & \emptycell & \emptycell \\
\SH & G & \NumRunsShuffle & \emptycell & \emptycell & \emptycell \\
\UN & G & --- & \emptycell & \emptycell & \emptycell \\
\bottomrule
\end{tabular}
\end{table}

\begin{table}[t]
\centering
\caption{H15's deciding cell (\TZ{} minus \FI, convention G, all leaks) and the cells beside it.
\pend}
\label{tab:h15}
\small
\begin{tabular}{lccccc}
\toprule
Cell & $b$ & $c$ & Difference (points) & Newcombe [$L$, $U$] & Labels \\
\midrule
Deciding: G, all leaks & \emptycell & \emptycell & \emptycell & \emptycell & \emptycell \\
Qualifier: M, all leaks & \emptycell & \emptycell & \emptycell & \emptycell & \emptycell \\
Qualifier: G, not constant & \emptycell & \emptycell & \emptycell & \emptycell & \emptycell \\
Placement-invariant subset & \emptycell & \emptycell & \emptycell & \emptycell & --- \\
Its complement & \emptycell & \emptycell & \emptycell & \emptycell & --- \\
Cut-date bootstrap & \multicolumn{4}{c}{\NumBootPLo th and \NumBootPHi th percentiles: \emptycell} & --- \\
\bottomrule
\end{tabular}
\end{table}

\paragraph{Registered sentences, to fill.} ``Truncation's rate lies between \pend{} and \pend{}
depending on how a missing bar is read (G to M: \pend{} to \pend).'' The read split of every
detector's catches \pend; the sensitivity ``ignores outputs that turn missing'' \pend; the
mirror asymmetry of each detector \pend; the controls against their predictions \pend; the union's
misses, classed by constancy under both definitions \pend; the comparisons at equal run counts
\pend; the descriptive run on seed \NumDescSeed{}, reported apart \pend.

% =============================================================================================
\section{What a verdict of the pass rule is worth}\label{sec:power}

Our forecasting comparisons (\cref{sec:case}) share one pass rule: a \uncited{Diebold--Mariano} test
on the hourly loss differences, with the \uncited{Harvey--Leybourne--Newbold} correction and a
rectangular kernel covering the overlap of the labels, \emph{and} a \uncited{stationary bootstrap}
interval of the mean difference on the same side of 0. We measured what a verdict under that rule is worth, on the
hourly series of the case study's first report, refitting nothing.

\paragraph{Size.} For each registered comparison, the centred hourly difference was sign-flipped
in circular blocks of $L \in \{\NumPowLa, \NumPowLb, \NumPowLc, \NumPowLd\}$ hours, with
\NumPowSizeReps{} replicates per $L$, and the rule was applied to each replicate as registered,
with the \uncited{Politis--White} block length recomputed per replicate (flip seed \NumPowSeed). The rule's
two-sided size was \NumPowSizeLo\pct{} to \NumPowSizeHi\pct{} for every comparison and $L$, none
above the \NumPowSizeCap\pct{} we had set as the alarm; in its registered direction,
\NumPowDirLo\pct{} to \NumPowDirHi\pct. The bootstrap alone reached \NumPowBootMax\pct{} (H10, at
$L=\NumPowBootLa$ and \NumPowBootLb), so the conjunction matters. The non-inferiority rule of H12
held \NumPowTwelveLo\pct{} to \NumPowTwelveHi\pct{} of the time when the true difference sat on its
margin.

\paragraph{Power.} Adding a shift $\delta$ in the registered direction to the flipped series, on a
grid of \NumPowGrid{} values with \NumPowGridReps{} replicates each, gives the smallest $\delta$
with \NumPowTarget\pct{} power, as a share of climatology's log-loss, for a constant shift and,
in parentheses, for a shift concentrated in \NumPowConcPct\pct{} of \NumPowConcBlock-hour blocks
($L=\NumPowConcL$): \texttt{matrix-1y} against \texttt{softmax-1y}, \NumPowSmLo\pct{} to
\NumPowSmHi\pct{} (\NumPowSmConc\pct); against \texttt{knn-1y}, \NumPowKnLo\pct{} to
\NumPowKnHi\pct{} (\NumPowKnConc\pct); H10, \NumPowTenLo\pct{} to \NumPowTenHi\pct{}
(\NumPowTenConc\pct); H11, \NumPowElLo\pct{} to \NumPowElHi\pct{} (\NumPowElConc\pct).
A constant shift is the easiest gain to see, so these are lower bounds on the detectable mean gain.

\paragraph{Reading.} The observed difference between \texttt{matrix-1y} and \texttt{knn-1y},
\NumPowKnObs\pct{} of climatology's log-loss, had \NumPowKnObsLo\pct{} to \NumPowKnObsHi\pct{}
power, so H3's ``does not hold'' against \texttt{knn-1y} is a difference too small for the rule to
see, not evidence that there is none. H11's difference, \NumPowElObs\pct, had \NumPowElObsLo\pct{}
to \NumPowElObsHi\pct{} power.
The study ran once (\NumPowSecs{}~s on \NumPowWorkers{} workers, \NumPowReps{} replicates in all);
a second run, meant as a byte-identity check, failed at the start-up of its worker pool before any
output, so this study has no identity check.

% =============================================================================================
\section{Matrix}\label{sec:matrix}

Matrix stores every entry $(z_i, y_i, t_i)$, a feature vector, an outcome and the time at which the
outcome became known, in time order, in an append-only log. Finished blocks of entries become leaves
of a binary tree whose nodes hold the sufficient statistics
\[
S = \Bigl(n,\ \textstyle\sum z,\ \sum z z^\top,\ \sum y z^\top,\ \sum y,\ \sum y^2\Bigr),
\]
merged by addition as a binary counter carries; a second tree holds the same nodes in centred form,
merged with the pairwise update of \uncited{Chan, Golub and LeVeque}. The raw entries are never
compacted away. A range of entries is covered by at most two partial blocks, scanned directly, and a
logarithmic number of nodes, so any window $R$ of entries yields its centred statistics, and its
ridge forecast for a query $z_q$ is computed at query time:
\[
A = M_{zz} + \Lambda,\qquad \beta = M_{yz} A^{-1},\qquad \hat y = \bar y + \beta\,(z_q - \bar z),
\]
with an unpenalised intercept and a penalty $\Lambda$ scaled to the fold's feature variances. The
solve is direct, never iterative; it is exact up to floating-point rounding, and it refuses a window
whose system lies within rounding of a singular one, by a bound we derived to first order and test
against an independent singular value decomposition but have not proved. The tree's nodes are sums
of entries, not a summary that replaces them.

Writing $a = A^{-1}(z_q - \bar z)$ in the centred coordinates, the forecast is a weighted sum of the
window's outcomes, so the contribution of any set of entries, down to a single day, is the same
product restricted to that set, and the contributions of a partition add up to the forecast; the
per-feature terms $\beta_j (z_{q,j} - \bar z_j)$ add up to the forecast minus the window's mean
outcome. This additivity is not
special to Matrix: any readout that is linear in stored outcome-side sums splits the same way,
Log-Linear Attention's mixture of bucket readouts included~\cite{guo2025loglinear}. % section 3.1, pp. 4-5
Causality is by construction: an entry joins Matrix only when its outcome is known, and a readout at
decision time $d$ over a window $w$ uses the entries with $d - w < t_i \le d$.

\paragraph{What we did not find.} Among the works we read, in the parts named in
\cref{sec:related}, none stores a raw per-entry log over which the ridge forecast of an arbitrary
window is computed directly from the window's sums. MesaNet solves a ridge regression over a gated
causal prefix by conjugate gradient~\cite{vonoswald2025mesanet}; % Eqs. (4), (6), (8); pp. 1-7
Log-Linear Attention keeps a Fenwick tree of linear-attention states over the prefix, with no
least-squares solve in the pages read~\cite{guo2025loglinear}; % pp. 1-8
Goulet Coulombe's identity relates a fixed training set to test rows, with no store and no
window~\cite{gouletcoulombe2025ols}. % pp. 1-10, A.1-A.6
For three further works the statement holds only in the parts we read, and some of those were read
through summaries: the regression cubes, whose cells aggregate losslessly over a lattice of stored
cells, with no arbitrary range, sliding window or subtraction of a cell in the pages
read~\cite{chen2006regcube}; % pp. 1-10 as images; pp. 11 on not read
the dual decomposition of Goulet Coulombe, G\"obel and Klieber, on a fixed training
set~\cite{gouletcoulombe2024dual}; % 2.1 (summarised)
and the least-squares modules of Garnelo and Czarnecki~\cite{garnelo2023}. % 2-3 (summarised)
The additive merge of regression statistics is 2006 prior art~\cite{chen2006regcube}, % Thms 4.3, 4.4
per-observation decomposition and negative weights are prior art~\cite{gouletcoulombe2024dual}, % 2.1, Table 2
and a causal prefix ridge with running statistics is prior art~\cite{vonoswald2025mesanet}. % Eq. (4), (6)
What remains is an engineering combination for an any-window, past-only memory; we claim no more.
\todo{OptMem, credited in our design notes, is not cited: read it first.}

% =============================================================================================
\section{Case study: pre-registered forecasts on hourly crypto bars}\label{sec:case}

\subsection{Setting}
Every hypothesis below was registered, with its decision rule, before its experiment ran, and is
reported here exactly as our results log records it. The panel holds the hourly bars of the members
of a monthly-rebalanced universe of Binance USDT pairs: \NumPanelMembers{} of the \NumPanelAll{}
members have a member month in the development period (the other \NumPanelLate{} join in 2025).
Rows exist only in a pair's member months; delisted pairs stay in the history. Each row has
\NumFeatures{} hand-made features and a three-class label (down, neutral, up) over
\NumLabelBars{} bars, entering at the next open. The test months are January 2019 to December 2024,
with monthly refits; the validation year 2025 and the vault are never read for a feature, a label or
a model. \NumScoredHours{} decision hours (\NumScoredRows{} rows, \NumScoredShare\pct{} of the
hours) are scored, the same for every arm; every statistic is computed on hourly means. Intervals are
\NumConf\pct{} stationary-bootstrap intervals. \NumPreviewMonths{} of the \NumTestMonths{} test
months (January 2019 to June 2020) were scored by smoke runs before the official run, with the same
arms and rules; no rule changed after them, and those months are also reported apart.

Matrix (\cref{sec:matrix}) is one pooled memory over all members; a readout uses the entries known
in the last \NumWinShort{} days, the last \NumWinYear{} days, or all history. The arms:
\code{climatology-panel} (class frequencies known at the fit); \code{matrix-select} (window and
penalty chosen at each fit on the prequential log-loss of rows whose labels are known); \code{matrix-1y}
(the \NumWinYear-day window); \code{softmax-1y} and \code{knn-1y}, a Nadaraya--Watson kernel and a
$k$-nearest-neighbour mean over the same entries, with their width or $k$ chosen at each fit;
\code{matrix-adaptive} (the penalty re-chosen at every decision); \code{matrix-stack} (three windows
stacked on the simplex); and \code{matrix-cusum} (refitted after a \uncited{Page CUSUM} alarm on
standardised recursive residuals, $k=\NumCusumK$, threshold \NumCusumH, average run length
\NumCusumARL{} days, instead of monthly). Every arm passed future invariance on a
\NumHonePanelMembers-member sub-panel at \NumHonePanelCuts{} cuts, with every honest output
identical to the bit and every planted online leak caught on its predicted pattern.

\subsection{\texorpdfstring{Report n\textsuperscript{o}\,2: Matrix}{Report 2: Matrix}}

\begin{table}[t]
\centering
\caption{Report n\textsuperscript{o}\,2, test months 2019--2024: mean hourly log-loss with
\NumConf\pct{} stationary-bootstrap intervals, and Brier skill score (BSS) against
\code{climatology-panel}; every BSS interval lies above 0. Variants tried so far:
\NumRepTwoVariants.}
\label{tab:report2}
\small
\begin{tabular}{lcc}
\toprule
Arm & Log-loss & BSS \\
\midrule
\code{climatology-panel} & \NumLLclim{} [\NumLLclimLo, \NumLLclimHi] & --- \\
\code{matrix-select} & \NumLLsel{} [\NumLLselLo, \NumLLselHi] & \NumBSSsel \\
\code{matrix-1y} & \NumLLoney{} [\NumLLoneyLo, \NumLLoneyHi] & \NumBSSoney \\
\code{softmax-1y} & \NumLLsoft{} [\NumLLsoftLo, \NumLLsoftHi] & \NumBSSsoft \\
\code{knn-1y} & \NumLLknn{} [\NumLLknnLo, \NumLLknnHi] & \NumBSSknn \\
\code{matrix-adaptive} & \NumLLadapt{} [\NumLLadaptLo, \NumLLadaptHi] & \NumBSSadapt \\
\code{matrix-stack} & \NumLLstack{} [\NumLLstackLo, \NumLLstackHi] & \NumBSSstack \\
\code{matrix-cusum} & \NumLLcusum{} [\NumLLcusumLo, \NumLLcusumHi] & \NumBSScusum \\
\bottomrule
\end{tabular}
\end{table}

A comparison of two arms \emph{passes} when the Diebold--Mariano $p$ is below \NumPassP{} and the
bootstrap interval of the difference lies on the same side of 0. Differences below are in mean hourly log-loss, negative when the first arm is
better, with the share of climatology's log-loss in parentheses.

\begin{description}[leftmargin=0pt, style=sameline]
\item[H2 holds.] \code{matrix-select} forecasts better than climatology: BSS \NumHtwoBSS{}
  [\NumHtwoBSSlo, \NumHtwoBSShi]. On the previewed months apart, \NumHtwoPrevBSS{}
  [\NumHtwoPrevLo, \NumHtwoPrevHi] over \NumHtwoPrevHours{} hours; on the other \NumRestMonths{}
  months, \NumHtwoRestBSS{} [\NumHtwoRestLo, \NumHtwoRestHi] over \NumHtwoRestHours{} hours.
\item[H3 does not hold.] H3 required \code{matrix-1y} to beat both \code{softmax-1y} and
  \code{knn-1y}, each comparison passing. Against \code{softmax-1y}: \NumHthreeSmDiff{}
  (\NumHthreeSmPct\pct), DM $p$ printed as \NumHthreeSmP, interval [\NumHthreeSmLo, \NumHthreeSmHi]:
  passes. Against \code{knn-1y}: \NumHthreeKnDiff{} (\NumHthreeKnPct\pct), $p$ \NumHthreeKnP,
  [\NumHthreeKnLo, \NumHthreeKnHi]: does not pass, with \NumPowKnObsLo--\NumPowKnObsHi\pct{} power
  at that size (\cref{sec:power}). \code{softmax-1y} is a fixed Nadaraya--Watson kernel on Matrix's
  entries, not a learned softmax attention; the passing comparison says nothing about softmax
  attention in general.
\item[H10 does not hold.] \code{matrix-adaptive} against \code{matrix-select}: \NumHtenDiff{}
  (\NumHtenPct\pct), $p$ \NumHtenP, [\NumHtenLo, \NumHtenHi]. The difference has the other sign: the
  adaptive arm's log-loss is the higher one. \NumHtenIdent\pct{} of hours are identical in the two
  arms.
\item[H11 does not hold.] \code{matrix-stack} against \code{matrix-select}: \NumHelevenDiff{}
  (\NumHelevenPct\pct), $p$ \NumHelevenP, [\NumHelevenLo, \NumHelevenHi], with
  \NumPowElObsLo--\NumPowElObsHi\pct{} power at that size.
\item[H12 holds.] \code{matrix-cusum} is non-inferior to monthly refits with at most half as many
  fits: difference \NumHtwelveDiff{} (\NumHtwelvePct\pct), one-sided \NumConf\pct{} upper bound
  \NumHtwelveUB{} and bootstrap \NumQuantPct th percentile \NumHtwelveQ, both below the margin of
  \NumMarginFactor{} times climatology's log-loss, \NumMargin; \NumHtwelveFits{} fits against a limit
  of \NumHtwelveLimit{} (\NumHtwelveAlarms{} alarms over \NumHtwelveDays{} daily values).
\end{description}

\subsection{\texorpdfstring{Report n\textsuperscript{o}\,3: an encoder with three attention operators}{Report 3: an encoder with three attention operators}}
A small encoder reads the \NumEncT{} hourly bars ending at the decision bar: \NumEncBlocks{}
pre-LayerNorm blocks of width \NumEncDmodel, \NumEncHeads{} heads of width \NumEncDhead{}, then a
Matrix over learned keys (a \NumEncMemDays-day window, entries joining when their label is known,
inside training episodes too) and a small decision head. Three arms differ only in the attention
operator. In \code{enc-ols}, each head computes $W = \tilde Q(\tilde K^\top \tilde K +
\Lambda)^{-1}\tilde K^\top$ and $O = WV$, with a constant key coordinate left unpenalised and a
penalty learned per head, so that every row of $W$ sums to one; \code{enc-softmax} is scaled
dot-product attention; \code{enc-linear} is \uncited{kernelised linear attention}. Parameters:
\NumParamsOls, \NumParamsSm{} and \NumParamsLin. Each arm is the mean of \NumSeeds{} seeds, with
\NumYearlyFits{} yearly fits per seed. All three arms passed future invariance on the sub-panel, bit
for bit, with every planted training-time leak on its predicted pattern; a positive control trained
on a label inside its own window reached a BSS of \NumEncPosBSS{} [\NumEncPosLo, \NumEncPosHi]. The
skill check of the planted controls failed on one of them, a memory stamped at decision time: H1
caught it to the bit, but its skill, \NumEncStampBSS{} [\NumEncStampLo, \NumEncStampHi], was below
the required \NumEncStampFloor. These are three verdicts, reported apart.

\begin{table}[t]
\centering
\caption{Report n\textsuperscript{o}\,3, test months 2019--2024: mean hourly log-loss and BSS with
\NumConf\pct{} intervals. Variants tried so far: \NumRepThreeVariants.}
\label{tab:report3}
\small
\begin{tabular}{lcc}
\toprule
Arm & Log-loss & BSS \\
\midrule
\code{enc-ols} & \NumLLols{} [\NumLLolsLo, \NumLLolsHi] & \NumBSSols{} [\NumBSSolsLo, \NumBSSolsHi] \\
\code{enc-softmax} & \NumLLencsm{} [\NumLLencsmLo, \NumLLencsmHi] & \NumBSSencsm{} [\NumBSSencsmLo, \NumBSSencsmHi] \\
\code{enc-linear} & \NumLLlin{} [\NumLLlinLo, \NumLLlinHi] & \NumBSSlin{} [\NumBSSlinLo, \NumBSSlinHi] \\
\bottomrule
\end{tabular}
\end{table}

\begin{description}[leftmargin=0pt, style=sameline]
\item[H4 holds, as non-inferiority.] Testable, since both arms have a BSS interval above 0.
  \code{enc-ols} minus \code{enc-softmax}: \NumHfourDiff{} (\NumHfourPct\pct); one-sided upper bound
  \NumHfourUB{} and bootstrap \NumQuantPct th percentile \NumHfourQ, both below the margin \NumHfourMargin. So
  \code{enc-ols} is no worse than \code{enc-softmax} by more than the margin. As a description, not
  a test: the bootstrap interval of the difference is [\NumHfourLo, \NumHfourHi] and the DM $p$ is
  \NumHfourP; no rule was registered for superiority, and we claim none.
\item[H5 does not hold.] \code{enc-ols} against \code{matrix-select}: \NumHfiveDiff{}
  (\NumHfivePct\pct), $p$ \NumHfiveP, [\NumHfiveLo, \NumHfiveHi]. No difference was shown in either
  direction. The two systems refit at different cadences, yearly and monthly, so H5 compares systems,
  not attention operators at equal cadence.
\item[H13: no verdict, the measurement is invalid.] H13 asked whether \code{enc-ols} is also cheaper
  to run, at most \NumHthirteenBound{} times \code{enc-softmax}'s median wall time over
  \NumHthirteenRepeats{} repeats. The ratio measured, \NumHthirteenRatio, met the rule's arithmetic,
  but it measured the timing stage's own GPU cache: its preparation left \NumCacheReserved{}~MiB
  reserved on a card of \NumCardTotal~MiB. On the same rows the ratio was \NumRatioAsIs{} as the code
  stands and \NumRatioEmptied{} with the cache emptied before each call, and synthetic rows give
  \NumRatioSynthLo{} to \NumRatioSynthHi{} on the same code path. The diagnosis was triggered by the
  gap with a synthetic reference on record before the run, not by the sign of the result. OLS-based attention is not shown to be cheaper to run than softmax
  attention here.
\item[Not run.] H6 (a Brier-loss ablation) was built and audited once but never fitted; H14 (a
  long/flat portfolio after costs) was registered and not run.
\end{description}

\subsection{Two exploratory analyses}\label{sec:case-explore}
Two readings of Report n\textsuperscript{o}\,2's skill were planned as exploratory, on test months
already seen, and run once on 2026-11-19 as a logged contact: the Brier and log-loss skill split into
``move or no move'' and ``direction given a move'', and a regression of the hourly loss differences on
the payoffs of reversal, momentum and volatility-timed momentum signals built from past bars, after
Nagel's reading of complex return forecasts~\cite{nagel2025}. % sections II.B-II.F, pp. 7-17
Because a payoff contains the realised return, that regression is an attribution after the fact,
not a test that the signals predicted the difference. Their choices were fixed before the call.
\todo{after audit 2a: both tables, labelled ``exploratory, test months''.}

% =============================================================================================
\section{Limits}\label{sec:limits}
\begin{itemize}[leftmargin=*]
\item \textbf{One generator.} Every rate is this generator's, on two pairs of daily bars and
  \NumCutDates{} cut dates. The population holds direct functions of the next bars; it holds no
  pipeline bug (a scaler fitted on the whole series, a centred window, a misaligned resampling), no
  fitted leak, no read of the time stamps and no read more than \NumReach{} bars ahead.
\item \textbf{Ties to future invariance.} Even under the quantile placement, the generator keeps ties to \FI's
  construction (\cref{sec:generator}), and natural thresholds fall inside \FI's bands at rates that
  differ by quantity (\cref{tab:bands}).
\item \textbf{Rebuilt rivals.} \TO{} rebuilds freqtrade's comparison on our expressions; it is not
  freqtrade's tool, and we do not trace how freqtrade loads its data. \RT{} equals windowed
  truncation on these functions. \TR's perturbation law is ours. None of these says what the
  original tools would do on real strategies.
\item \textbf{Fitted models.} Truncation and re-tiling, as run here, apply to fit-free functions
  only.
\item \textbf{The platform.} The population of a seed depends on the platform's exponential; the
  results hold for the pinned environment. The code digest checked by the runs pins the Windows
  working tree, with its line endings, rather than the commit, whose identity is its hash.
\item \textbf{Tests deferred.} Five mutants planted by the last code audit survive the detector
  tests; the frozen code is right on all five by the audit's reading, and their tests come after the
  runs.
\item \textbf{The vault.} The validation year (2025) and the vault (2026 onwards) are sealed in code,
  and the vault months were moved out of the data folder in a first physical stage, in plaintext, on
  2026-10-29. In an audit of the bench on 2026-10-28, two tests written to prove holes in the code
  seal read
  BTCUSDT daily vault rows outside the seal: the \NumExposedBars{} bars of the March 2026 raw archive
  and the vault months of the stored file. They assert row counts, and one prints the first bar of
  that archive when it succeeds; no feature, model or decision used these rows. The exposure is
  stated wherever the vault is used. \todo{The protocol's split table gives the vault as 2026-01 to
  2026-10, while the data notes count stored bars to 2026-09: settle the wording. The exposure is
  dated 2026-10-28, after Audit A of iteration 2, in the protocol's change log; the paper's brief
  calls it iteration 1's: settle which.}
\item \textbf{The case study.} \NumPreviewMonths{} of its \NumTestMonths{} test months were seen
  through smoke runs before the official run. Its exploratory analyses run on months already seen.
  H5 compares different refit cadences; H13 has no valid measurement, and the GPU wall times of
  Report n\textsuperscript{o}\,3's cost table carry the same artefact. Future invariance of the
  encoder arms tested the pipeline of their first seed, while each arm's model is the mean of its
  \NumSeeds{} seeds.
\item \textbf{Who audited.} The adversarial audits were made by an AI model distinct from the one
  that wrote the code (\cref{sec:ai}), not by an independent human reviewer.
\item \textbf{The reading.} Every source was read in part, most through a summarising fetch tool
  (\cref{sec:related}).
\end{itemize}

% =============================================================================================
\section{Reproducibility}\label{sec:repro}
\todo{Repository URL and archive (Abdou's decision).} The code is MIT-licensed; the data are not in
the repository (\cref{sec:data}). Every run writes its configuration, its commit, a digest of the
code and a digest of the data manifest.

\paragraph{Commits and runs.} Report n\textsuperscript{o}\,2 ran at \code{31c8005} on Windows
(\path{runs/2026-10-29-report2/}); its re-run with the hourly series at \code{221591c}
(\path{runs/2026-10-30-report2/}), and in WSL2 at \code{5f17a76} (\path{runs/2026-10-30-report2-wsl/}).
Report n\textsuperscript{o}\,3 ran at \code{5f17a76} in WSL2 on an RTX~3060 GPU
(\path{runs/2026-10-30-report3/}); the power study at \code{5f17a76} (\path{runs/2026-10-30-power/}).
The iteration-3 adversary ran at \code{3a22daa} and the future-shuffle twins at \code{221591c}. The
benchmark's code is frozen at \code{b71ba65}; its runs are \path{runs/2026-11-19-detectors-descriptive/},
\path{runs/2026-11-19-detectors-confirmatory/} and its rerun, made at \code{f12b9b5}. The case
study's exploratory call ran at \code{1088aac} (\path{runs/2026-11-19-casestudy/}).

\paragraph{\texorpdfstring{\code{-{}-check}}{--check} identities.} Re-run with \code{-{}-check}, Report n\textsuperscript{o}\,2
reproduces its \NumRepTwoCheckFiles{} files byte for byte (\code{run.json}, which holds the wall
clock, excluded), and its re-run its \NumRepTwoRerunFiles{} files, the hourly series included;
Report n\textsuperscript{o}\,3 reproduces its \NumRepThreeCheckFiles{} files. The iteration-3
adversary and the future-shuffle twins were each run twice, with byte-identical outputs. The
benchmark's confirmatory run was repeated with another number of workers: \pend.

\paragraph{Two platforms.} Windows~11 (CPython \NumPyWin) and Ubuntu in WSL2 (CPython \NumPyWSL)
differ in the last bits of \code{np.log} and \code{np.exp}. Report n\textsuperscript{o}\,2 re-run in
WSL2 kept every scored row, label, time and symbol; probabilities differed in \NumPlatProbArms{} of
\NumPlatArms{} arms by at most \NumPlatProbMax, and \NumPlatNums{} of the report's \NumPlatNumsOf{}
numbers differed, by at most \NumPlatRelMax{} relative; every verdict was identical. For the leak
populations the same last-bit difference changes which leaks a seed draws, which is why the benchmark
pins its platform (\cref{sec:prereg}).

% =============================================================================================
\section{Use of AI tools}\label{sec:ai}
This project was built with Claude Code agents (Anthropic). An orchestrating session planned each
iteration and delegated the code, the experiments, the literature reading and the adversarial audits
to agents with fixed roles; the audits were made by a model other than the one that wrote the code
under audit. The literature was read through a web-fetch tool that summarises pages, which is why
each citation names the part read (\cref{sec:related}). This first draft was written by such an agent
from the project's records, and every number in it is generated from those records by a script.
The author, \todo{name}, takes full responsibility for every claim of this paper. AI tools are not authors.

% =============================================================================================
\section{Data and licence}\label{sec:data}
\textbf{Source: Binance Vision (data.binance.vision), CC BY-NC-SA 4.0.} The dataset terms that govern
these archives~\cite{binancevisionterms} % clauses 3, 4.1, 4.5 (read 2026-10-25; re-read 2026-11-18 through a summarising tool)
license them under CC BY-NC-SA 4.0, allow academic research, and require a redistributed derivative
work to credit Binance Vision and to be licensed under CC BY-NC-SA 4.0. This paper derives figures
and tables from those archives; nothing in it implies that Binance sponsors or endorses it. The
data themselves are not redistributed. Our reading of the terms, and of the Binance Terms of Use
that they supplement, is partial and is not legal advice. \todo{Licence of this paper: Abdou's
decision, pending; CC BY-NC-SA 4.0 is what the dataset terms' clause 4.5 imposes on derived works.}
The leak corpus's code is ours and MIT-licensed; its outputs on the bars derive from the data, and
their release waits for an answer on the terms. \todo{Release of the corpus: Abdou's decision.}

\bibliographystyle{plainnat}
\bibliography{refs}

\end{document}
